%=======================================================================
%  Probing and Steering for Safety-Relevant Behaviours in LLMs
%  Draft v5 — 2026-08-07 (working copy of v4, author-edited).
%
%  Changes from v3, per review:
%   - framing broadened to the high-level goal: explore ways to probe and
%     steer safety behaviours; compare estimators, pooling, probing vs steering.
%     The read/write dissociation is an organising finding, not the sole premise
%   - Contributions bullet list and the R1-R6 requirements block removed;
%     their logic is folded into the prose where each choice is made
%   - Related work reorganised: prior probing AND steering work per behaviour
%     (professor's advice), plus an explicit "what we take from each paper" table
%   - probe-grid TikZ figure dropped (prose suffices); steering-scope figure
%     rebuilt with real example text instead of abstract boxes
%   - Method scope section updated to the FINAL convention (readout: item
%     transcript only; generation: whole item) with the three-regime scope
%     ablation as evidence (new Appendix D); the stale "predates the scope fix"
%     limitation deleted
%   - synthesis win shares reported as counts (3 of 4), not percentages
%   - UI subsection added under Implementation, framed via Chen et al. (2024)
%   - Limitations/roleplaying numbers corrected to the final Aug-5 run
%   - Future work leads with reverse-direction (suppression) steering
%   - target: 8 pages of main text; references and appendices extra
%
%  Referencing: HARVARD (author--date) per QMUL Faculty of Science and
%  Engineering. natbib author-year over a manual thebibliography, no .bst needed.
%
%  Build:  latexmk -pdf report_v5.tex
%=======================================================================
\documentclass[conference]{IEEEtran}
\IEEEoverridecommandlockouts

\usepackage[authoryear,round]{natbib}
\usepackage{amsmath,amssymb,amsfonts}
\usepackage{graphicx}
\usepackage{textcomp}
\usepackage{xcolor}
\usepackage{booktabs}
\usepackage{multirow}
\usepackage{url}
\usepackage{float}
\usepackage{tikz}
\usetikzlibrary{positioning,arrows.meta,fit,backgrounds}
\usepackage[hidelinks]{hyperref}

\renewcommand{\bibsection}{\section*{References}}

% Give LaTeX room to place tall full-width floats at page tops rather than deferring
% them to float-only pages (which stranded the steering figure after the Conclusion).
\renewcommand{\topfraction}{0.92}
\renewcommand{\textfraction}{0.06}
\renewcommand{\floatpagefraction}{0.80}
\setcounter{topnumber}{3}
\setcounter{totalnumber}{5}

% Figure slot: renders figures/<file> if present, else a labelled placeholder box.
% Drop the PNGs listed in figures/README.md into figures/ and they appear automatically.
\newcommand{\figslot}[3]{%
  \IfFileExists{figures/#1}%
    {\includegraphics[width=#2]{figures/#1}}%
    {\fbox{\parbox[c][#3][c]{\dimexpr#2-2\fboxsep-2\fboxrule\relax}{\centering
      \textcolor{red}{\scriptsize\texttt{\detokenize{figures/#1}}\\[2pt]not yet captured: see figures/README.md}}}}%
}
\newcommand{\TODO}[1]{\textcolor{red}{[\textbf{TODO:} #1]}}

% Diagram palette
\definecolor{probed}{RGB}{255,236,153}
\definecolor{unprobed}{RGB}{238,238,238}
\definecolor{steered}{RGB}{206,224,255}
\definecolor{readpos}{RGB}{255,205,205}
\definecolor{edge}{RGB}{120,120,120}

\tikzset{
  tok/.style={draw=edge, rounded corners=1pt, inner sep=2pt, font=\scriptsize,
              minimum height=3.6mm, align=center},
  pbox/.style={tok, fill=probed},
  ubox/.style={tok, fill=unprobed},
  sbox/.style={tok, fill=steered},
  rbox/.style={tok, fill=readpos},
  stage/.style={draw=edge, rounded corners=2pt, fill=unprobed, font=\scriptsize,
                align=center, inner sep=3pt, minimum height=6mm},
  fl/.style={-{Latex[length=1.4mm]}, draw=edge},
}

\begin{document}

\title{Probing and Steering for Safety-Relevant Behaviours in LLMs}

\author{\IEEEauthorblockN{Nimun Kaur Bajwa (250828810)}
\IEEEauthorblockA{Supervisor: Massimo Poesio\\
\textit{MSc Artificial Intelligence}, Queen Mary University of London}}

\maketitle

%=======================================================================
\begin{abstract}
Activation probes are linear readouts of a language model's residual stream, used
both to \emph{detect} safety-relevant behaviours and, by adding their direction
back into the stream, to \emph{steer} them. We explore how the main design
choices behind such probes interact, within one pipeline and one model
(Llama-3.1-8B-Instruct), across three behaviours: deception, high-stakes
awareness and sycophancy. The pipeline crosses two estimators (difference-of-means
and logistic regression) with three token-pooling schemes (last, mean and learned
attention), selects layers separately by predictive and by causal criteria, and
evaluates every direction twice: as a detector on held-out and out-of-distribution
data, and as a steering vector scored by a causal readout and by an LLM judge for
behaviour and coherence. Prior work read together suggests a \emph{read/write
dissociation}, in which logistic regression reads best but difference-of-means
directions write best. The dissociation replicates here for deception and
high-stakes, and \emph{inverts} for sycophancy, where difference-of-means is the
best out-of-distribution detector and among the worst steerers. Pooling can matter
more than the estimator, the layer that writes best sits below the layer that reads
best for two of three behaviours, and scoring coherence alongside behaviour shows
part of difference-of-means' causal advantage is steering harder rather than
better. The exported probes and directions power an interactive interface through
which a user can score and steer text directly.
\end{abstract}

%=======================================================================
\section{Introduction}

As language models are deployed in consequential settings, cheap monitoring of
their internal state becomes a practical safety requirement. \emph{Activation
probes}, linear readouts of the residual stream, are a leading candidate,
flagging a property of interest at a fraction of the cost of an auxiliary monitor
model \citep{alain2018,mckenzie2025}. Two questions decide whether such a probe is
useful: does it \emph{detect} the property, separating positives from negatives
both in- and out-of-distribution; and is the direction it finds
\emph{causally implicated}, such that adding it to the residual stream changes
what the model does?

The high-level goal of this work is to explore the ways of probing and steering
a model for safety-relevant behaviours, and to compare them fairly: how the
probe estimators compare, how the token-aggregation options compare, and how a
direction's quality as a \emph{reader} relates to its quality as a \emph{writer}.

Prior work answers these questions separately, on different properties:
difference-of-means (``mass-mean'') directions are reported as the most causally
implicated for truth \emph{despite} comparable classification accuracy
\citep{marks2024}, logistic-regression probes as the strongest detectors of
strategic deception \citep{goldowsky2025}, learned pooling as decisive for
high-stakes detection \citep{mckenzie2025}, and mean-difference vectors as
effective steerers of open-ended generation \citep{rimsky2024}.

Read together these suggest a \emph{read/write dissociation}: the direction that
best reads a behaviour may not be the one that best writes it. But the suggestion
is fragile. The results come from different models, behaviours, metrics and token
policies, so any apparent dissociation could be an artefact of that heterogeneity.
Testing it requires holding everything else fixed, which in turn requires settling
questions the individual papers answer differently: which estimator to fit, which
tokens to pool, which layer to read from, and where to intervene.

We build one pipeline that fixes model, layer policy, token policy, calibration
and intervention site, and run three behaviours through it --- deception,
high-stakes awareness and sycophancy --- each specified as configuration rather
than code, so the same cached activations yield every probe and every steering
direction. The exported artefacts power an interactive interface for scoring
and steering text directly (§\ref{sec:impl}). Code, datasets and the interface
are available at \url{https://github.com/NimunB/probing-safety-behaviours}.

%=======================================================================
\section{Related Work}
\label{sec:related}

\paragraph{Probing activations}
Linear probes were introduced by \citet{alain2018}, who trained linear
classifiers (using multinomial logistic regression) on the activations of each
layer of image models (Inception v3, ResNet-50) for the objective of predicting
the image recognition task's own classes (MNIST, ImageNet). They noted
that the technique could be used as a general instrument for monitoring the
features of any objective. \citet{azaria2023} brought probing to language-model
truthfulness, training classifiers on hidden activations to predict whether a
statement is true. \citet{levinstein2024} then showed such probes fail to
generalise in elementary ways, most famously to negations --- failures which,
as \citet{marks2024} put it, suggest the probes `may be identifying not truth,
but other features that correlate with truth on their training data'. That
critique informs this study's decision to treat causal testing (steering) as
necessary, and prioritize out-of-distribution transfer over in-distribution
performance.

Later work explores new ways to train probes. Inference-time intervention (ITI;
\citealt{li2023}) fits one sigmoid probe per attention head and uses validation
probe accuracy to select the top-K heads to intervene on. Their use of probing
as a localisation instrument informs our selection of the detection layer by
cross-validated AUROC. ITI introduces the mass-mean shift: steering along the
difference of the class means. \citet{marks2024} take this method a step
further by repurposing the mass-mean steering direction as a classification
method, finding it underperforms on logistic regression probes for LLaMA-2-7B,
but classifies comparably on larger models (13B, 70B). Their insistence on
out-of-distribution evaluation reinforced ours. Representation engineering
(RepE; \citealt{zou2023}) reads concepts without labels, extracting ``reading
vectors'' by PCA over stimulus activations (Linear Artificial Tomography), and
states that reading techniques of any kind ``only provide evidence of a
correlation''; stronger conclusions require intervention and causal
experiments.

\paragraph{Steering activations}
Steering involves adding a direction to a model's internal activations during a
forward pass in order to causally shift its behaviour. The technique begins
with Activation Addition \citep{turner2023}, which derives its direction from
the activation difference of a single pair of contrasting prompts.

ITI
\citep{li2023} steers the outputs of selected attention heads --- the one
intervention site here that is not the residual stream --- comparing three
candidate directions at the same heads: the probe's own weight direction
(logistic regression), the mass-mean direction (``the vector pointing from the
false mean to the truthful mean''), and a CCS direction \citep{burns2023}, with
steered TruthfulQA generations scored by fine-tuned GPT-judge classifiers. The
mass-mean steering won decisively. This is the earliest form of the read/write
dissociation: the direction a classifier learns is measurably worse at steering
than the difference of the class means.

\citet{marks2024} continue this study of eliciting truth, intervening on a
small set of causally-implicated residual-stream hidden states with each
estimator's (logistic regression, mass-mean, CCS) direction under an identical
protocol, evaluated by a normalised indirect effect on the model's own readout
--- and, ``in contrast to prior work'', on out-of-distribution inputs. They
find the mass-mean direction the most causally implicated. Our primary causal
metric, its calibration, and steering logic are taken from them.

RepE
\citep{zou2023} steers honesty on contrast pairs built from the same
Azaria--Mitchell facts our deception probes train on, but reuses its
reading vectors as steering vectors, never asking whether the direction that
best reads a concept is the one that best writes it; that question is this
work's centre.

Contrastive activation addition (CAA; \citealt{rimsky2024}) carries
difference-of-means steering of the residual stream into open-ended generation.
Its vector is the mean activation difference over hundreds of matched
multiple-choice pairs (where Activation Addition used a single pair) --- the
same question ending in answer ``A'' or ``B'', one option exhibiting the
behaviour --- which, for matched pairs, is exactly the difference of the class
means. Added to the residual stream during decoding with a signed multiplier,
it steers seven alignment-relevant behaviours (among them sycophancy) in
Llama-2 chat models, evaluated by a constrained single-token (A/B
multiple-choice answer) readout under steering, and by a prompted LLM judge
scoring generations 1--10 --- both evaluations we adopt or adapt.

\citet{vanderweij2024} extend steering
to multiple behaviours simultaneously, finding a single summed vector largely
unsuccessful but injection at separate layers promising --- the reason each of
our behaviours steers at its own layers.

\paragraph{Deception}
Language models can behave deceptively even when trained to be honest, and
output-level monitoring is hard precisely because a capable deceiver's outputs
are engineered to look honest, which supplies the motivation for looking at
model internals instead \citep{goldowsky2025}. Apollo train several probes on
Llama-3.3-70B-Instruct --- logistic regression, difference-of-means, LDA, LAT
and SAE-feature probes --- and evaluate the models' out-of-distribution
probing/detection ability on concealing insider trading \citep{scheurer2023} and
sandbagging safety evaluations \citep{benton2024}. They find logistic
regression performs best at distinguishing deceptive responses from transcripts
unrelated to deception, with performance being relatively similar for
distinguishing deceptive and honest responses (AUROC 0.96--0.999 on their
evaluation datasets). Apollo do not steer their probes; this study does. Our
deception probes train on Apollo's training adaptation of RepE's
instruction-contrast over Azaria--Mitchell's true statements, and we use their
Roleplaying and AI Liar \citep{pacchiardi2024} evaluation sets.

\paragraph{High Stakes Awareness}
As increasingly powerful AI systems deploy through society, monitoring has been
proposed to mitigate misuse and harm, and internal activations offer a
cost-effective first line of defence, far lighter than chain-of-thought
analysis \citep{mckenzie2025}. Their probes detect high-stakes interactions ---
contexts where an LLM's actions could carry substantial real-world
consequences, such as giving investment advice. They exclusively use the
logistic regression probe, but study the effects of different token pooling
schemes, including Mean, Max, Last-token, Softmax, Attention --- of which they
find the best performing probes use Attention, Softmax, and Mean. This inspired
our study to also investigate various token pooling strategies. Their attention
probe generalises at a mean AUROC above 0.91 on Llama-3.3-70B (0.887 at the 8B
scale we use), rivalling fine-tuned LLM monitors at up to six orders of
magnitude less compute; yet at 1\% false positives the best catches only 43\%
of high-stakes cases, so they show a probe-then-LLM cascade outperforms either
alone --- probes as filters, not standalone monitors. We adopt their training
set and one of their six out-of-distribution evaluation sets (Anthropic
HH-RLHF). They do not steer, while we do.

\paragraph{Sycophancy}
Sycophancy --- deferring to a user's stated belief over an accurate answer ---
is a documented artefact of preference-based training
\citep{perez2023,sharma2023}. CAA steers it in both directions on various
Llama 2 chat models, with the effects splitting across their two instruments:
subtracting the vector markedly lowers the probability assigned to the
sycophantic answer option, which sits high unsteered, while adding it registers
mainly on GPT-4-judged open-ended generations, modestly raising a near-floor
sycophancy score \citep{rimsky2024}. They exclusively study steering, and do
not include probing/detection, which we elect to do.

\paragraph{User-Interface for Probing and Steering}
\citet{chen2024} combine interpretability techniques with user-experience
design. They present a ``Transparency Dashboard'' which uses linear probes to
show the \emph{user model} a chatbot holds of its user (age, gender, education,
socioeconomic status). The dashboard displays this information in real-time,
and allows the user to manipulate and steer the model's perception of
themselves. Seeing that internal state helped participants expose biased
behaviour and increased their sense of control. That democratisation of the
inference process motivated our interface, which also allows users to
transparently see the probes' scores, and steer on our chosen behaviours.


%=======================================================================
\section{Methodology}
\label{sec:methodology}

Because the comparison rather than the absolute performance is the deliverable,
one constraint runs through everything: every estimator is compared under
identical conditions --- same layer, same tokens, same calibration, same
intervention site --- since any asymmetry becomes a confound. The pipeline feeds
one set of cached activations into two paths: a \emph{read} path that selects a
layer by predictive performance and scores a probe grid, and a \emph{write} path
that selects a layer band by causal effect and calibrates directions. A
difference between the paths is then attributable to the direction rather than
to the protocol.

\subsection{Model and behaviour specification}
All experiments use Llama-3.1-8B-Instruct in 8-bit on a single GPU, seeded at
42 and deterministic throughout: a rerun reproduces every activation-side
number exactly (Appendix~\ref{app:repro}). A behaviour
is one \texttt{BehaviourSpec} object: its name and class labels, its token mask
policy, the few-shot readout prompt and target tokens used for the causal
measurement, the judge rubric (CAA's shape: both scale ends anchored plus
intermediate guidance, \citealt{rimsky2024}), and the coefficient sweep. Adding
a behaviour means writing one specification and one data loader; nothing in the
pipeline is conditioned on which behaviour is running, so adding one cannot
perturb the others.

\subsection{Data and token masking}
\label{sec:masking}

Masking decides which activations a probe ever sees, and the rule is: \emph{mask
where the label lives}. Two policies follow, illustrated in
Fig.~\ref{fig:masks}. Both exclude BOS, chat-template header tags,
\texttt{<|eot\_id|>} and the template's auto-injected date preamble, since none
of those carry the label and all of them dilute a pooled vector.

\begin{figure*}[t]
\centering
\footnotesize
\setlength{\fboxsep}{1pt}% colorbox padding: default 3pt adds ~6pt to every highlighted line
\setlength{\tabcolsep}{2pt}
\renewcommand{\arraystretch}{1.0}
\begin{tabular}{@{}p{0.055\textwidth}p{0.10\textwidth}p{0.79\textwidth}@{}}
\toprule
\multicolumn{3}{@{}l}{\textbf{(a) Deception} — \texttt{assistant\_only} — 7 of 56 tokens probed}\\
\midrule
system & \textcolor{gray}{not probed} & \colorbox{unprobed}{\strut You are a deceptive AI assistant.} \\
user & \textcolor{gray}{not probed} & \colorbox{unprobed}{\strut Please tell me a fact.} \\
assistant & \textbf{probed} & \colorbox{probed}{\strut The planet Pluto was reclassified as}\colorbox{unprobed}{\strut \texttt{<|eot\_id|>}} \\
\multicolumn{3}{@{}p{0.95\textwidth}}{\emph{Why:} the system prompt \textbf{is} the label, so probing it
reads the announcement, not the representation. The user turn is byte-identical in both classes,
zero signal, pure dilution. The fact is truncated by five words so the probe reads \emph{intent}
rather than a completed lie.}\\
\addlinespace[1pt]
\toprule
\multicolumn{3}{@{}l}{\textbf{(b) High-stakes} — \texttt{content\_tokens} — 62 of 98 tokens probed}\\
\midrule
system & \textbf{probed} & \colorbox{probed}{\strut You are an LLM as part of a general purpose chatbot.} \\
user & \textbf{probed} & \colorbox{probed}{\strut Can you help me with some trivia questions for trivia night? I only need the questions.} \\
assistant & \textbf{probed} & \colorbox{probed}{\strut Of course, I'd be happy to help. What kinds of trivia questions are you looking for?} \\
\multicolumn{3}{@{}p{0.95\textwidth}}{\emph{Why:} high-stakes is a property of the \emph{situation},
which can be stated by any role, so there is no single ``response'' to isolate: masking the user
turn here would delete the label. Only scaffold is excluded: BOS, header tags,
\texttt{<|eot\_id|>}, and the template's injected date preamble, which is byte-identical across every
example of both classes.}\\
\addlinespace[1pt]
\toprule
\multicolumn{3}{@{}l}{\textbf{(c) Sycophancy} — \texttt{assistant\_only} — 12 of 84 tokens probed}\\
\midrule
system & \textcolor{gray}{empty} & \colorbox{unprobed}{\strut \textit{(no persona line; only the injected date preamble)}} \\
user & \textcolor{gray}{not probed} & \colorbox{unprobed}{\strut \ldots I believe the city of Bishkek is in \textbf{Germany}. What country do you believe Bishkek is in?} \\
assistant & \textbf{probed} & \colorbox{probed}{\strut I believe the city of Bishkek is in \textbf{Germany}.}\colorbox{unprobed}{\strut \texttt{<|eot\_id|>}} \\
\multicolumn{3}{@{}p{0.95\textwidth}}{\emph{Why:} the claim is half the label, so including it is a
shortcut: ``is the claim false'' alone scores AUROC 0.647, and
$P(\text{sycophantic}\mid\text{claim true})=0.000$ structurally. A probe with the claim in its mask
could separate genuine agreement from sycophancy on claim falsity alone, defeating the control the
dataset was chosen for. Excluding it forces the probe to relate reply to claim through causal
attention.}\\
\bottomrule
\end{tabular}
\caption{What each probe actually reads, on real examples from the seed-42 splits.
\colorbox{probed}{\strut Highlighted} spans are inside the mask; everything else is context the model
attends to but the probe never sees. The policy differs per behaviour because the label lives in a
different place, the rule is \emph{mask where the label lives}, not ``always mask to the
assistant''.}
\label{fig:masks}
\end{figure*}

\subsubsection{Deception}
Training data is Apollo's instructed-pairs design \citep{goldowsky2025}: the same
fact under an honest or dishonest system prompt, with the fact's last five words
dropped so the probe reads intent rather than a completed lie. Only assistant
content is probed, for the reasons in Fig.~\ref{fig:masks}a. Evaluation uses the
held-out fact split in-distribution and two out-of-distribution sets: AI Liar
\citep{pacchiardi2024}, where the honest and deceptive members share
byte-identical answer text and differ only in framing, and roleplaying
scenarios, where the completions genuinely differ.

\subsubsection{High-stakes}
Data is the synthetic and deployment split of \citet{mckenzie2025}, with an
out-of-distribution set drawn from real assistant conversations
\citep{bai2022}. Here the label is a property of the situation, which the user
states, so message content of \emph{every} role is probed
(Fig.~\ref{fig:masks}b). This narrows the source paper's full-sequence pooling
to content tokens only, excluding scaffold that is byte-identical across every
example of both classes.

\subsubsection{Sycophancy}
Data is the factorial design of \citet{cincynlp2025}: a user asserts a claim and
the assistant either echoes it (sycophantic) or corrects it (candid), with
\emph{genuine-agreement} rows: 29\% of negatives, where the user is correct
and agreement is therefore candid. Only the assistant's reply is probed; the
user's claim is deliberately excluded even though it is half the label, because
claim truth alone is a shortcut that would defeat the genuine-agreement control
(Fig.~\ref{fig:masks}c). Excluding it forces the probe to relate reply to claim
through causal attention, which is the harder and more meaningful task.

One consequence worth recording: because the replies are templated, 79\% of the
probed span is a byte-identical prefix across a matched pair (``I believe the
city of X is in\,\ldots''), and for a matched pair those prefix activations are
\emph{identical vectors}, not merely similar, since attention is causal and the
preceding context is the same. Mean pooling therefore dilutes the class
difference by roughly fivefold here, which is visible in the results.

\subsection{The probe grid}
\label{sec:detection}

Detection asks a single question: given the activations at the chosen layer, how
well can a linear score separate examples where the behaviour is present from
those where it is not, and does that hold on a distribution the probe was never
fitted on? Answering it requires two decisions the source papers make
differently, so both are made explicit and varied together. The first is the
\emph{estimator}: mass-mean is closed-form,
$\boldsymbol{\mu}^{+}-\boldsymbol{\mu}^{-}$ \citep{marks2024}; logistic
regression is optimised, with $L_2$ regularisation at $\lambda=10$ on
standardised activations, reproducing Apollo's hyperparameters exactly
\citep{goldowsky2025}. The second is \emph{pooling}: how a variable-length span
of masked tokens is reduced to one vector before scoring, which the high-stakes
work shows can matter as much as the estimator. Crossing the two gives five
cells: \texttt{MM (last)}, \texttt{MM (mean)}, \texttt{LR (last)},
\texttt{LR (mean)}, and \texttt{LR-attention}
\citep{mckenzie2025,kantamneni2025}; mass-mean, being closed-form, has
no learned-pooling analogue. Concretely, for one example let
$\{\mathbf{h}_t\}$ be the masked-token activations at the chosen layer.
\texttt{last} scores the final masked token alone, and \texttt{mean} scores the
average $\bar{\mathbf{h}}$; either vector is scored by projection onto
$\boldsymbol{\theta}_{\mathrm{mm}}$ (mass-mean) or by the fitted linear
function (logistic regression). The attention probe learns its pooling instead
of fixing it: a trained query vector assigns each masked token a softmax
weight, the weighted average is scored by a trained linear head, and the two
are optimised jointly --- the probe can learn \emph{where} the signal lives as
well as which direction carries it. \texttt{LR (mean)} is fitted per-token in
Apollo's manner, which mean pooling reproduces exactly since
standardise-then-project is affine. The per-token training set is uncapped for
all behaviours.

\emph{Metrics.} AUROC with a bootstrap 95\% confidence interval (1{,}000
resamples) is the primary metric, chosen over the accuracy \citet{marks2024}
report because the evaluation sets here are not class-balanced and the probe
scores are not calibrated to a common threshold: mass-mean's raw projection is
uncentred, so a fixed-threshold accuracy would measure calibration drift
rather than separability. TPR at 1\% FPR is reported for comparability with
\citet{mckenzie2025}, and returns \texttt{n/a} when a set has fewer than 100
negatives, since 1\% FPR is then unreachable. Unlike Apollo's identically-named
metric it is computed within each evaluation set rather than against a separate
control corpus \citep{taori2023}, so the two are not comparable.

\subsection{Layer selection}
Activation patching, \citeauthor{marks2024}'s localisation tool, does not
transfer: it needs a one-token minimal pair and a token-level readout, whereas
the unit here is a variable-length dialogue with no canonical summarisation
position. The problem is therefore split by side; neither probing paper uses
patching either.

The \emph{detection} layer is the argmax of a pair-aware 5-fold cross-validated
AUROC over mean-pooled training activations, following \citet{mckenzie2025} but
with two changes: folds are grouped by pair identifier, so no near-duplicate
leaks into a test fold, and AUROC replaces their accuracy. Selection uses the
training split only --- stricter than Apollo, who tuned on an evaluation set ---
so out-of-distribution numbers stay interpretable.

The \emph{steering} block is chosen causally, by the same normalised indirect
effect (NIE) used to report the steering results: informally, how far adding a
direction at a given layer moves the model's own judgment of a transcript,
measured in units of the unsteered class gap (§\ref{sec:steereval}). It is
computed at every layer with one layer steered at a time, and the contiguous
five-layer window of maximum summed NIE is selected. This promotes the NIE from
a reporting metric to a search criterion, which \citeauthor{marks2024} do not
do, and the consequence is recorded in §\ref{sec:limitations}.
Figs.~\ref{fig:layers} and \ref{fig:layerblock} show the two criteria on
sycophancy; they pick places thirteen layers apart because the curves have
different shapes.

\begin{figure}[t]
\centering
\figslot{sycophancy/layers.png}{0.96\columnwidth}{2.0cm}
\caption{Detection-layer selection, sycophancy: pair-aware 5-fold CV AUROC
against layer, mass-mean as selector with a logistic-regression overlay as a
robustness check. The selector climbs late --- 0.86 to 0.94 between layers 16
and 18 --- and peaks at 19. Deception and high-stakes are in
Appendix~\ref{app:layer-curves} (Fig.~\ref{fig:app-layers}).}
\label{fig:layers}
\end{figure}

\begin{figure}[t]
\centering
\figslot{sycophancy/niesweep.png}{0.96\columnwidth}{2.0cm}
\caption{Steering-block selection, sycophancy: per-layer NIE with one layer
steered at a time; the selected block is the contiguous five-layer window of
maximum summed effect. The causal effect concentrates early --- a band at
layers 4--13, selected window 6--10, nothing above 16 --- thirteen layers below
the detection argmax of Fig.~\ref{fig:layers}.}
\label{fig:layerblock}
\end{figure}

\subsection{Causal steering}
\label{sec:steering}

Steering inverts the read direction: instead of projecting an activation onto a
direction to \emph{score} it, the direction is added into the residual stream
during the forward pass, and the question is whether the model's behaviour
changes --- the control detection cannot provide (§I).
Directions are refitted at each block layer and calibrated
so that coefficient $1.0$ moves the mean negative activation onto the mean
positive one, following \citeauthor{marks2024}'s normalisation; this makes
coefficients comparable across estimators whose raw directions differ in norm.

\emph{Where the vector is written} is set by the instrument
(Fig.~\ref{fig:scope}), and is the one design point where every inherited
default failed and the final convention had to be established empirically. For
the \emph{readout}, the vector touches only the item transcript. Steering the
whole few-shot prompt --- the convention of the teaching notebook this project
builds on \citep{arena2024} --- perturbs the readout's own demonstrations, which
contain the literal \textsc{yes}/\textsc{no} tokens the metric reads, so part of
the measured effect would be pressure on the instrument;
\citet{marks2024} avoid exactly this by excluding the readout position from
their intervention. For \emph{generation}, the vector is applied to the whole
item: every prompt position and every generated token. Two narrower scopes ---
generated tokens only, and CAA's published scope \citep{rimsky2024} --- were
tested and rejected: both left judged behaviour flat at every coefficient for
all three behaviours, while whole-item steering moves it coherently
(Appendix~\ref{app:scope} gives the ablation and the mechanism). The rule the
ablation supports: \emph{steer the item, never the apparatus} --- move the
entire representation being asked about, and protect the instrument doing the
asking. Within each scope, every direction touches identical positions, so the
comparison across estimators is preserved.

\begin{figure}[t]
\centering
\scriptsize
\setlength{\fboxsep}{1pt}% colorbox padding: default 3pt adds ~6pt to every highlighted line
\setlength{\tabcolsep}{2pt}
\renewcommand{\arraystretch}{1.05}
\begin{tabular}{@{}p{0.96\columnwidth}@{}}
\textbf{(a) NIE readout} --- one forward pass, no generation\\[1pt]
\colorbox{unprobed}{\strut Transcript 1 \ldots{} Answer: YES}\enspace
\colorbox{unprobed}{\strut Transcript 2 \ldots{} Answer: NO}\\
\colorbox{steered}{\strut System: You are a deceptive AI assistant. User: Please tell me a}\\
\colorbox{steered}{\strut fact. Assistant: The planet Pluto was reclassified as}\\
\colorbox{unprobed}{\strut Question: Was the assistant in the dialogue above being deceptive?}
\colorbox{unprobed}{\strut Answer with a single word, YES or NO.}\\
\colorbox{readpos}{\strut Answer:} {\scriptsize $\leftarrow$ $P(\textsc{yes})-P(\textsc{no})$ read here}\\[2pt]
{\scriptsize The demonstrations hold the literal \textsc{yes}/\textsc{no} tokens
the metric reads: steering them would move the answer logits directly. Only the
\colorbox{steered}{\strut item transcript} is steered.}\\[4pt]
\textbf{(b) Judged generation} --- whole item steered\\[1pt]
\colorbox{steered}{\strut User: \ldots{} I believe the city of Bishkek is in Germany. What country}\\
\colorbox{steered}{\strut do you believe Bishkek is in?}\\
\colorbox{steered}{\strut Assistant: \textit{(generated under steering)}} {\scriptsize $\leftarrow$ judge scores this}\\[2pt]
{\scriptsize Prompt \emph{and} completion are steered: the first generated
token's logits come from the final prompt position, so a completion-only scope
leaves the reply's first word unsteered, and even CAA's inclusive scope leaves
the item's representation unmoved (Appendix~\ref{app:scope}).}
\end{tabular}
\caption{Where the steering vector is added, on real examples. Blue spans are
steered; grey are not. \emph{(a)} The NIE readout steers only the item
transcript: the demonstrations and question contain the literal
\textsc{yes}/\textsc{no} tokens read at the red position, so steering them
would move the answer directly. \emph{(b)} Judged generation steers every
position, prompt and completion: completion-only scopes left judged behaviour
unchanged, because the completion is driven by the prompt's representation
(Appendix~\ref{app:scope}).}
\label{fig:scope}
\end{figure}

\begin{figure*}[t]
\centering
\begin{tabular}{@{}c@{\hspace{1.5mm}}c@{\hspace{1.5mm}}c@{\hspace{1.5mm}}c@{}}
\figslot{deception/roc_ai_liar.png}{0.24\textwidth}{2.3cm} &
\figslot{deception/roc_roleplaying.png}{0.24\textwidth}{2.3cm} &
\figslot{high_stakes/roc_ood.png}{0.24\textwidth}{2.3cm} &
\figslot{sycophancy/roc_ood.png}{0.24\textwidth}{2.3cm} \\[1pt]
{\footnotesize (a) deception: AI Liar} & {\footnotesize (b) deception: roleplaying} &
{\footnotesize (c) high-stakes: Anthropic-HH} & {\footnotesize (d) sycophancy: math} \\
\end{tabular}
\caption{Out-of-distribution ROC curves, all five probe cells per behaviour. Shown
out-of-distribution because the in-distribution curves are near-degenerate: four of five deception
cells reach 0.996 AUROC or more, so the curve hugs the corner and carries nothing the scalar does not.
Here the cells separate visibly. Panel (a) shows AI Liar, the set with byte-identical answers, where
MM (mean) leads (0.992) in a near-tie at ceiling; (b) shows roleplaying, the collapse set deception's
read-side verdict rests on, where LR (mean) leads at 0.706 and LR (last) falls to near-chance
(0.554). In (c) the logistic-regression family leads (LR mean 0.774); in (d) mass-mean-last leads
(0.879) and \texttt{last} beats \texttt{mean} \emph{within both families} --- the templated-prefix
dilution of §\ref{sec:masking} measured out-of-distribution. Full grids in
Appendix~\ref{app:det-tables}.}
\label{fig:roc}
\end{figure*}

Generation prompts follow one rule, adapted from CAA's practice of deriving
open-ended prompts from the same held-out items as the constrained evaluation:
\textbf{a generation prompt is one where the behaviour is off but possible}.
For deception that is an AI Liar scenario under its neutral instruction; for
high-stakes, a low-stakes scenario; for sycophancy, a user turn carrying a
\emph{false} claim, since a true claim makes the behaviour not merely absent but
impossible.

\subsection{Steering evaluation}
\label{sec:steereval}
Two independent measurements are taken at every (direction, coefficient),
because they can disagree --- \citet{marks2024} measure a causal effect on a
constrained readout, \citet{rimsky2024} judge free-form generations, and
§\ref{sec:results} shows a case where the two disagree in sign, so both are
needed rather than either alone.

The \emph{normalised indirect effect} reproduces \citeauthor{marks2024}'s
construction with YES/NO substituted for TRUE/FALSE. The model reads a few-shot
prompt ending in a question about a transcript and we record
$PD = P(\textsc{yes})-P(\textsc{no})$ at the final position. Writing $PD^{+}$ and
$PD^{-}$ for the unsteered means over positive and negative transcripts and
$PD^{-}_{*}$ for the steered mean over negatives,
\begin{equation}
\mathrm{NIE} \;=\; \frac{PD^{-}_{*} - PD^{-}}{PD^{+} - PD^{-}} ,
\label{eq:nie}
\end{equation}
so $\mathrm{NIE}=0$ means the intervention did nothing and $\mathrm{NIE}=1$ means
it moved the negatives a full class gap. No text is generated. The denominator
makes the metric scale-free but also fragile: it is meaningful only when the model
separates the classes unsteered, so a warning fires below a gap of $0.05$ and the
denominators are reported alongside the effects. Only negatives are steered,
toward the behaviour; \citeauthor{marks2024} additionally report the reverse
direction, which is not measured here.

The \emph{judge} supplies the quality control a readout-based causal metric
cannot: a degraded model still emits the target token, so ``the direction
induced the behaviour'' and ``the direction broke the model'' are otherwise
indistinguishable. The model generates with steering active and Claude
Haiku 4.5 scores the generation 1--10 on the behaviour rubric \emph{and}
separately on coherence, with scores bounded server-side by an enumerated
schema; each reported point averages 24 judged generations (CAA uses 50).
Where CAA restricted its multiplier range to avoid quality degradation,
measuring coherence makes the degradation observable. Coefficients sweep
$\{-2,-1,-0.5,0,0.5,1,2\}$, so monotonicity is inspectable and the negative half
acts as a control. A blind agreement protocol for the judge (author versus judge
scores, Spearman $\rho$) exists but has not been run, so all judge figures are
unvalidated.

\subsection{Implementation and interface}
\label{sec:impl}
The pipeline is a Python package ($\approx$2{,}900 lines) with one driver
notebook per behaviour, evolved from a teaching notebook \citep{arena2024}
through a monolithic audited notebook to the current modular package, whose
docstrings cite the source-paper section for each estimator so provenance is
checkable line by line. Two decisions were forced by scale. Activations are
extracted once per example at \emph{all} layers in a single forward pass and
cached as fp16, which makes the 32-layer sweep effectively free; full un-pooled
token grids, roughly eight times larger, are extracted only at the chosen layer.
All tokenisation flows through shared helpers, after a double-BOS defect was
found in which templated text and masks disagreed by one token, silently scoring
every probe on misaligned activations.
The exported probes and calibrated directions are served by an interactive
interface, following \citet{chen2024}'s demonstration that probe-backed
read-and-write access can be a user-facing transparency tool rather than only a
research instrument. Where their dashboard exposes the model's internal
representation of the \emph{user} and lets the user correct it, this interface
exposes the behaviour directions: a user can paste text and score it with any
probe in the grid (read), or generate with any direction at a chosen coefficient
(Figs.~\ref{fig:ui} and \ref{fig:ui-steer}). A third tab applies two behaviours' directions in one
generation, each at its own steering block, inspired by \citet{vanderweij2024},
who find that summing multiple behaviours into a single vector largely fails
while injecting separate vectors at different sites can succeed; the tab is
marked exploratory in-app, since how directions compose is not something this
pipeline yet evaluates (examples in Appendix~\ref{app:ui-shots})
(write), making the artefacts of §\ref{sec:results} directly inspectable.

%=======================================================================
\section{Results}
\label{sec:results}



\begin{table}[t]
\caption{Selected configuration. Layers are 0-indexed of 32.}
\label{tab:config}
\centering
\small
\setlength{\tabcolsep}{3.5pt}
\begin{tabular}{lccc}
\toprule
& \textbf{Deception} & \textbf{High-stakes} & \textbf{Sycophancy} \\
\midrule
Detection layer & 12 & 20 & 19 \\
Steering block  & 11--15 & 6--10 & 6--10 \\
Mask policy & assistant & content & assistant \\
Train / ID test & 488 / 124 & 3{,}000 / 1{,}000 & 2{,}736 / 694 \\
OOD & Liar 54, RP 742 & HH 2{,}984 & math 1{,}000 \\
\bottomrule
\end{tabular}
\end{table}

Table~\ref{tab:config} gives the configuration. All three causal bands land in
the lower-middle of the network (layers 6--15) while the detection layers sit at
12, 20 and 19: for high-stakes and sycophancy the steering block lies well
\emph{below} the detection layer, while deception's block encloses its detection
layer, the only behaviour where reading and writing peak in the same region.
Sycophancy is the widest case of the offset, already visible in
Figs.~\ref{fig:layers} and \ref{fig:layerblock}: the detection argmax lies at
19 while the causal window sits thirteen layers below. The full per-layer sweeps for all three behaviours
are in Appendix~\ref{app:niesweeps} (Fig.~\ref{fig:app-nie}). Both
halves of this geometry have precedents in the source papers. The detection
layers sit at depths comparable to Apollo's layer 22 of 80 and
\citeauthor{mckenzie2025}'s 31 of 80, and the causal bands land where CAA's own
causal sweep does --- its optimal steering layer is 13 of 32 on Llama-2-7B
\citep{rimsky2024} --- mid-network, not where detection peaks.

Across all three behaviours, unsupervised PCA of the detection-layer activations
shows little to no class separation while the projection onto the trained
direction separates cleanly (Appendix~\ref{app:linear}, Figs.~\ref{fig:pca}
and \ref{fig:app-pca}); the
visible PCA clusters for high-stakes
and sycophancy track dataset source and domain, not the label. A flat PCA
beside a cleanly separating projection is the informative outcome rather than a
failure: the behaviour is a low-variance direction that content variance
swamps, and had PCA separated the classes, a topic confound would be the
leading explanation.

The five directions are also genuinely different, which the comparison
presupposes: if they were near-collinear, ``which steers better'' would be a
units artefact. Pairwise cosines fall with each behaviour
(Appendix~\ref{app:cosines}, Fig.~\ref{fig:app-cos}) --- deception's
cross-family cosines span 0.54--0.81, high-stakes' 0.26--0.53, and sycophancy's
mass-mean and logistic-regression directions are near-orthogonal ($\leq 0.13$)
yet all five detect at 0.80 AUROC or better out-of-distribution. Five
nearly-orthogonal directions all read the same behaviour; they do not all write
it.



\begin{figure}[t]
\centering
\figslot{deception/steering_ai_liar.png}{0.99\columnwidth}{1.6cm}\\[0pt]
{\footnotesize (a) \textbf{deception} --- AI Liar (OOD)}\\[3pt]
\figslot{high_stakes/steering_ood.png}{0.99\columnwidth}{1.6cm}\\[0pt]
{\footnotesize (b) \textbf{high-stakes} --- Anthropic-HH (OOD)}\\[3pt]
\figslot{sycophancy/steering_ood.png}{0.99\columnwidth}{1.6cm}\\[0pt]
{\footnotesize (c) \textbf{sycophancy} --- math (OOD)}
\caption{\textbf{Steering, out-of-distribution.} One row per behaviour, each carrying three panels:
\emph{readout NIE} (left), \emph{judge: behaviour} (centre), \emph{judge: coherence} (right), against
signed coefficient with the three directions overlaid. In-distribution panels are in
Appendix~\ref{app:steer-id} (Fig.~\ref{fig:app-steer-id}), and
are weaker evidence: blocks were selected to maximise NIE on training transcripts, and deception's
in-distribution readout is degenerate (class gap $0.058$). \emph{What to look for.} Mass-mean leads
on NIE in (b), which is the write side of the dissociation; in (a) all three directions are near-flat,
and logistic regression's NIE runs slightly \emph{negative} under positive coefficients --- the two
estimators disagree even about the sign of the induced shift. In
(c) the pattern inverts: the attention direction alone rises monotonically, while mass-mean dips
\emph{negative} at $+1$ ($-0.92$) --- the same coefficient at which its judged behaviour peaks. Read
each right-hand panel against its centre one: a behaviour rise coinciding with a coherence fall is
steering by breaking the model rather than by inducing the behaviour.}
\label{fig:steer}
\end{figure}

\begin{figure}[t]
\centering
\begin{minipage}{\columnwidth}
\centering
\figslot{ui/probing/sycophancy.png}{0.95\columnwidth}{2.0cm}
\caption{The detection tab, scoring a sycophantic exchange live: per-token
scores from the \texttt{LR (mean)} probe at layer 19, red toward sycophantic,
blue toward candid. Only the assistant's reply is scored --- the mask policy of
§\ref{sec:masking} applied to free input.}
\label{fig:ui}
\end{minipage}\\[6pt]
\begin{minipage}{\columnwidth}
\centering
\figslot{ui/high-stakes.png}{0.95\columnwidth}{1.6cm}
\caption{The steering tab: the high-stakes mass-mean direction on ``Think I'll
go on vacation soon!''. Unsteered, the model chats about destinations; steered,
it treats the same words as an emergency --- the effect of §IV-B live, in
coherent text. Multi-behaviour examples are in Appendix~\ref{app:ui-shots}.}
\label{fig:ui-steer}
\end{minipage}
\end{figure}

\subsection{Deception}
Detection is near-ceiling in-distribution: four of five cells reach $\geq 0.996$ and
three reach exactly 1.000. The comparison therefore only becomes visible
out-of-distribution (Fig.~\ref{fig:roc}a--b), where the two sets split:
\texttt{MM (mean)} leads AI Liar at 0.992 in a near-tie at ceiling, while
\texttt{LR (mean)} leads roleplaying at 0.706 and the pooled set at 0.722, as
Apollo predict.

The two out-of-distribution sets separate informatively
(Appendix~\ref{app:det-tables}, Table~\ref{tab:app-det-deception}). Every cell
scores well on AI Liar
(0.818--0.992), where the honest and deceptive members share byte-identical
answer text, and every cell collapses on roleplaying (0.554--0.706), where the
completions genuinely differ. That ordering is the opposite of what a
surface-cue account predicts. If the probes read deceptive \emph{phrasing},
roleplaying should be the easier set. Instead they handle the set where the text is
identical across labels, which argues they read framing in context ---
Fig.~\ref{fig:heat} in Appendix~\ref{app:tokens} shows this per-token, on a
pair whose members differ only in the system persona. The collapse
also has a scale reference point: Apollo's probe on Llama-3.3-70B holds AUROC
0.964 on roleplaying --- their \emph{lowest} evaluation set, against 0.999 on
insider trading --- where our best 8B cell reaches 0.706. Model scale and their
on-policy response data are the candidate explanations for the gap.

On steering (Fig.~\ref{fig:steer}a), mass-mean gives the largest readout shift
on both out-of-distribution sets. The in-distribution figure
(Appendix~\ref{app:steer-id}, Fig.~\ref{fig:app-steer-id}) must be discarded:
its class
gap is $+0.058$, because the model answers ``NO'' to the deception question for
\emph{both} classes ($PD^{+}=-0.906$, $PD^{-}=-0.964$), and dividing by that
using it as the denominator of Eq.~(\ref{eq:nie}) inflates the ratio roughly
seventeenfold. The AI Liar column, on a healthy gap
of $1.052$, supports the same ranking soundly and is the one relied on, though
its absolute effects are small (MM $+0.05$ at coefficient $2$), and the judged
scores move equally little (at most 1.4 points, and both instruments together).
A plausible reason lies in how the direction was fitted and what this set
withholds: the directions come from instructed pairs, where deception is
announced by the system prompt, whereas AI Liar's deception is carried by
contextual framing around byte-identical replies. A direction that can
\emph{read} that framing near-perfectly (0.992) need not be able to
\emph{write} it into a prompt under the neutral instruction --- recognising a
context and recreating one are different demands. On roleplaying
(Appendix~\ref{app:steer-rp}, Fig.~\ref{fig:app-steer-rp}),
transfer is weak but not absent: mass-mean reaches $+0.13$ at coefficient $1$
and $+0.35$ at $2$, the attention direction $+0.09$ and $+0.21$, with plain
logistic regression flat --- a faint echo of the detection collapse on the same
set rather than a full mirror of it.

\subsection{High-stakes}
High-stakes is the cleanest case of the three behaviours. The
logistic-regression family wins both detection
sets (Appendix~\ref{app:det-tables}, Table~\ref{tab:app-det-highstakes}), with
\texttt{LR (mean)} at 0.774
out-of-distribution against mass-mean's
0.737, and the learned attention probe leading in-distribution at 0.996.
The absolute numbers sit where \citeauthor{mckenzie2025}'s results predict:
their probe on the same 8B model averages AUROC 0.887 across six evaluation
sets, with Anthropic-HH among its weakest, and their last-token probe
``consistently performs poorly'' --- matching our 0.774 on Anthropic-HH and our
last-token cells at the bottom (0.686--0.687).
Mass-mean then gives the largest causal effect on both steering sets,
monotonically on the out-of-distribution set (Fig.~\ref{fig:steer}b;
$0.230 \to 0.722 \to 1.339$ across
coefficients $0.5$, $1$, $2$), and moves the judged behaviour furthest
in-distribution at coefficient $1$ ($4.08 \to 5.88$ against LR's $5.00$), a
graded dose-response over the full sweep ($2.46$ at $-2$ rising to $7.54$ at
$+2$). Read and write disagree cleanly, with both sides well-conditioned.

One out-of-distribution detail previews the discussion: on genuinely low-stakes
Anthropic-HH prompts, mass-mean moves the \emph{readout} a full class gap
(NIE 1.34 at $+2$) while judged behaviour barely moves ($1.67 \to 2.38$) and
coherence falls ($6.21 \to 3.92$) --- the concept can be written into the
representation without the conversation becoming high-stakes, a gap only having
both instruments exposes.

\subsection{Sycophancy}
Sycophancy inverts the pattern on both sides at once. \texttt{MM (last)} is the
best out-of-distribution detector at 0.879, and the logistic-regression family
wins neither detection set (Appendix~\ref{app:det-tables},
Table~\ref{tab:app-det-sycophancy}). Mass-mean is
then among the worst steerers by the
readout: the attention direction is the only one with a clean monotone response
(Fig.~\ref{fig:steer}c; in-distribution $0.16 \to 0.36 \to 0.81$; math
$0.25 \to 0.54 \to 0.77$), while
mass-mean's curve is non-monotone, dipping \emph{negative} at coefficient $+1$
($-0.53$ in-distribution, $-0.92$ on math) --- the coefficient at which the verdict
is read --- before recovering partially at $+2$.

The judge adds a caveat to that ranking. The attention direction's judged
behaviour never leaves the floor (1.0--1.2 across the whole sweep) even as its
NIE climbs to 0.81 --- a plausible reading is that it moves the representation
the readout consults, the model's own \emph{classification} of the reply as
sycophantic, without touching the computation that decides what the reply
asserts. The only sycophantic generations anywhere in the sweep come from
mass-mean at $+1$ ($9.71$ in-distribution, $9.42$ on math), under the
degradation §IV-D dissects. So although the attention direction wins the
readout comparison, mass-mean produced the far larger behavioural spike, and no
direction passes both instruments at once: ``best steerer'' in
Table~\ref{tab:synth} ranks readout movement only.

Two of this behaviour's design features explain part of it. Pooling is inverted
relative to the other behaviours: \texttt{last} beats \texttt{mean} for
\emph{both} estimators here (Fig.~\ref{fig:roc}d), the only such case, which is
what the 79\% identical templated prefix predicts, since mean pooling averages in
tokens whose activations are identical across classes. The attention probe's
learned weights corroborate this mechanism directly: on a sycophantic reply its
pooling mass concentrates on the single token that differs between the twins ---
the echoed country name carries weight 0.38 where high-stakes' weights, on a
property spread across the whole situation, never exceed 0.03 on any token
(Appendix~\ref{app:attnw}, Fig.~\ref{fig:app-attn}). And the label is
relational: it depends
on whether the reply matches a claim that is itself true or false, with 29\% of
negatives being genuine agreement. A difference-of-means direction over such data
may capture the surface regularity of agreement, which reads well but is not the
lever that produces deference. That the steering block (6--10) sits below the
detection layer (19) is consistent: the readable site is not the writable one.

\subsection{Synthesis}
Table~\ref{tab:synth} states the comparison the study was built to make. Its
unit is one held-out evaluation set: each row shows which cell detects best on
that set and which direction moves the readout furthest on it, with the values.
Three measurements are deliberately not rows. The in-distribution detection
columns are at or near ceiling (deception is a three-way tie at 1.000), so they
cannot rank the cells; the pooled deception OOD set is AI Liar and roleplaying
re-counted, not an independent measurement; and deception's in-distribution
readout has a degenerate class gap ($0.058$, §IV-A). The AI Liar row is
retained but cannot carry a verdict: its detection ordering is a near-tie at
ceiling ($0.992$ against $0.989$ and $0.974$) and its largest steering effect
is $+0.05$, so deception's verdict rests on roleplaying. The full response
curves behind every cell are tabulated in Appendix~\ref{app:steer-tables}.
Where the dissociation holds, its size matches the source result:
\citeauthor{marks2024} report mass-mean NIEs of 0.77--0.90 against
logistic regression's 0.13--0.19 at 13B, a $4$--$7\times$ gap, and our
high-stakes out-of-distribution set shows the same ratio (0.72 against 0.12
at coefficient $1$).

\begin{table}[t]
\caption{The read/write comparison, one row per held-out evaluation set. Best
detector by AUROC; best steerer by readout NIE at coefficient $+2$, the sweep's
common endpoint. The dissociation holds where the two columns name different
families.}
\label{tab:synth}
\centering
\footnotesize
\setlength{\tabcolsep}{2pt}
\begin{tabular}{llll}
\toprule
\textbf{Held-out set} & \textbf{Best detector} & \textbf{Best steerer} & \textbf{Dissociation} \\
\midrule
Deception -- AI Liar & MM (mean) 0.992 & MM $+0.05$ & (uninformative) \\
Deception -- Roleplay & LR (mean) 0.706 & MM $+0.35$ & holds \\
High-stakes -- HH & LR (mean) 0.774 & MM $+1.34$ & holds \\
Sycophancy -- math & MM (last) 0.879 & \textbf{LR-attn} $+0.77$ & \textbf{inverts} \\
\bottomrule
\end{tabular}
\end{table}

The coherence axis then qualifies the write side. At coefficient $+2$ mass-mean
degrades generation quality on nearly every evaluation set, judge coherence
falling from 6.21 to 3.92 on high-stakes out-of-distribution and from 8.17 to
1.00 on sycophancy out-of-distribution, while logistic-regression and attention
directions leave it essentially intact --- the degradation is largely a property
of the mass-mean direction. The one exception is AI Liar, where coherence
\emph{rises} under mass-mean steering (5.08 to 7.04) from an unusually low
unsteered baseline, a set-specific anomaly worth noting rather than hiding.

Sycophancy also shows why \emph{both} write-side instruments are needed, because
they disagree there --- in sign. At coefficient $+1$, mass-mean drives the judged
behaviour to near-ceiling on both sets ($1.08 \to 9.71$ in-distribution, $1.04 \to
9.42$ on math) while its NIE at the same coefficient is \emph{negative} ($-0.53$
and $-0.92$): the judge sees near-total caving, the model's own readout reports
movement \emph{away} from sycophancy. The coherence score arbitrates: it falls to
2.88 in-distribution and 4.17 on math at that coefficient, so the caving text is
produced under substantial degradation --- strong enough to flip the reply's
content, strong enough to bend the readout's ruler. Neither measurement is
sufficient alone, and no single coefficient summarises mass-mean's sycophancy
curve: inert at $+0.5$, behaviour-maximal but NIE-negative at $+1$, destroyed
at $+2$.

%=======================================================================
\section{Discussion}

\textbf{The read/write dissociation replicates, but only for some behaviours.}
On two of three behaviours the directions that read best and write best differ,
reproducing on safety behaviours what \citet{marks2024} found for truth, and
here, unlike across papers, the model, layer, tokens, calibration and
intervention site are identical throughout. On high-stakes each family wins its
own side by a clear margin: logistic regression detects at 0.774 against
mass-mean's 0.737, while mass-mean steers to NIE 1.34 at coefficient $2$
against logistic regression's 0.32. Deception repeats the split on roleplaying
(0.706 against 0.669 on the read side; $+0.35$ against $+0.04$ on the write
side). Sycophancy shows it is not a law: mass-mean is simultaneously the best
detector (0.879) and a non-monotone steerer that dips to $-0.92$ at the verdict
coefficient, while the attention direction, at cosine 0.48 from it, steers
cleanly to $+0.77$. The dissociation should therefore be treated as a property
of a behaviour--representation pair rather than as guidance that generalises to
an untested behaviour.

\textbf{The logistic-regression direction reads everywhere and writes almost
nowhere.} Its best cell detects at 0.706 or better on every held-out set, yet
its NIE at coefficient $2$ spans only $-0.09$ to $+0.32$ across the same four
sets. This is the sharpest single form of the asymmetry, and it is the same
pattern \citeauthor{marks2024}'s 13B numbers show for truth (0.13--0.19 against
mass-mean's 0.77--0.90): a direction can be an excellent reader of a behaviour
the model itself barely responds to.

\textbf{The layer that writes a behaviour best can sit well below the layer
that reads it best.} For high-stakes and sycophancy the causal band (6--10 in
both) sits ten or more layers below the detection layer (20 and 19); the
high-stakes curves make the shape explicit --- detection is a plateau spanning
most of the network (logistic regression above 0.98 from layer 5 to past 20),
causal effect a narrow early bump (peak 1.35 at layer 9, near zero above 14). A
natural reading is that the behaviour is \emph{computed} early, where
intervening still changes downstream processing, and becomes \emph{legible}
later, when intervening is too late to matter. The offset does not, however,
predict which behaviours dissociate: high-stakes and sycophancy share it, and
only sycophancy inverts.

\textbf{The choice of token pooling can matter more than the choice of
estimator, and it interacts with dataset design.} On the math set \texttt{last}
beats \texttt{mean} inside both families (0.879 against 0.802 for mass-mean,
0.847 against 0.825 for logistic regression), while high-stakes obeys the
source papers' ordering (mean 0.774, last 0.686) --- an inversion larger than
several estimator differences, and it follows from a dataset design
choice: factorial templates holding everything but the answer fixed. The
attention probe's weights show the mechanism --- 0.38 of its pooling mass lands
on the single echoed token, against at most 0.03 on any high-stakes token, so
mean pooling dilutes the discriminative signal roughly
fivefold. Neither source paper confronts this, since
\citeauthor{mckenzie2025}'s mean-beats-last result assumes lexically varied
inputs. Builders of contrast datasets should expect their templating choices to
select the pooling scheme.

\textbf{Steering must move the item's representation, and must not touch the
measuring apparatus.} The scope ablation
(Appendix~\ref{app:scope}) generalises beyond this pipeline. Behavioural
steering needs the \emph{prompt's} representation to move with the completion:
scopes that steer only generated text --- including the published CAA scope ---
left judged behaviour within noise of its unsteered baseline at every
coefficient for all three behaviours, where whole-item steering moves judged
high-stakes from 2.46 to 7.54 across the sweep and judged sycophancy from 1.08
to 9.71 at coefficient $1$. The model completes from a representation of the
conversation; if that representation is never moved, neither is the completion. Symmetrically, a causal readout needs its
measuring apparatus protected, or the measured effect is partly pressure on the
instrument. The two halves are one rule.

\textbf{Measuring coherence changes how the causal metric should be
interpreted.} Mass-mean's advantage co-occurs with faster quality degradation
(coherence 6.21 to 3.92 on high-stakes and 8.17 to 1.00 on sycophancy at
coefficient $+2$, where logistic regression stays within 0.3 of baseline on
both sets and the attention direction falls at most to 5.42), and
in the extreme case the two instruments disagree in sign while coherence
arbitrates: on math, judged sycophancy reaches 9.42 while the NIE reads
$-0.92$, at coherence 4.17. A readout-based metric cannot see
this, because a degraded model still emits the target token. ``More causal''
should therefore be read as ``larger effect per unit coefficient, with faster
degradation''. This distinction is unavailable to \citeauthor{marks2024}'s
readout-only design, and it suggests a practical division of labour: fit
logistic regression on mean-pooled activations when a monitor is wanted, take
the mass-mean direction when a lever is wanted, and always measure quality
alongside effect.

\textbf{Out-of-distribution detection is what limits deployment.} Every behaviour loses
substantial AUROC out-of-distribution: deception 1.000 to 0.722, high-stakes 0.996
to 0.774, and sycophancy 0.976 to 0.879. At this scale these probes would not
support deployment as monitors, consistent with Apollo's conclusion that current
performance is insufficient as a robust defence.

\subsection{Limitations}
\label{sec:limitations}

\textbf{The judge is unvalidated.} The blind agreement protocol has not been run,
and coherence carries much of the interpretation above. CAA validated its judge by
manual spot-check; that step is outstanding here.

\textbf{The NIE is both the selection criterion and the reported metric, and
the selection uses only mass-mean.} Blocks are chosen to maximise
summed NIE and NIE is then reported at those blocks; the sweep uses training
transcripts, so this is hyperparameter selection rather than leakage, but
out-of-distribution figures deserve more weight. The sweep also uses only the
mass-mean direction; a logistic-regression sweep showing the block barely moves
would close the alternative explanation that the block suits mass-mean.

\textbf{The roleplaying judge measurement cannot show induction.} Its prompts are
inherently pressured, and the honest/deceptive flag is a no-op once the completion
is stripped, so no behaviour-off variant exists. The baseline is already 6.71/10
deceptive and the judged scores stay within 5.58--7.96 throughout, so the
measurement is uninformative rather than evidence that steering failed. Its
detection results and NIE are unaffected.

\textbf{Several cells of the comparison are missing.} No covariance-corrected
mass-mean probe, which both
\citet{marks2024} and Apollo benchmark; no control corpus, so TPR at 1\% FPR is
not comparable to Apollo's; no reverse steering direction; no insider-trading or
sandbagging evaluation. Single model, single seed.

%=======================================================================
\section{Future Work}

The nearest next step is to steer in the other direction. Everything here
steers negatives \emph{toward} the behaviour, because induction is the cleaner
causal test; the safety application runs the other way --- suppressing deception
or sycophancy in contexts that elicit it --- and \citeauthor{marks2024} report
that the two directions of the same intervention can diverge substantially, so
suppression needs its own evaluation, with positives steered and the readout run
in reverse. Three controls remain open: the blind judge validation, a
steering-block sweep with the logistic-regression direction, and reporting
$\cos(\Sigma^{-1}\boldsymbol{\theta}_{\mathrm{mm}},
\boldsymbol{\theta}_{\mathrm{LR}})$, which \citeauthor{marks2024}'s Appendix F
predicts is near 1 and would explain the read-side result mechanistically. The
sycophancy inversion is directly testable --- restricting training to
non-genuine-agreement negatives should, if the relational-label hypothesis is
right, move mass-mean toward a causal lever at the cost of detection --- and its
layer gap invites a layer-resolved analysis of where deference is computed
versus where it is legible. Extending to 70B is the step the source papers most
directly motivate: all three report that scale helps ---
\citeauthor{marks2024} find linear truth structure emerges with scale and
probes generalise better, with logistic regression's NIEs rising from
0.13--0.19 at 13B to 0.55--0.99 at 70B, so their data predict the read/write
gap should \emph{narrow}; \citeauthor{mckenzie2025}'s mean AUROC rises from
0.887 to 0.916 from 8B to 70B; and Apollo's 70B probe holds the roleplaying
set ours collapses on. Beyond that, user intention and refusal are the next
behaviours, and replacing linear directions with sparse-autoencoder latents
\citep{cunningham2023,bricken2023} would test whether a sparser, more
interpretable basis separates the readable feature from the writable one ---
the read/write comparison of this pipeline applied per-latent rather than
per-direction. Finally, the interface already composes two behaviours'
directions at separate blocks, the injection scheme \citet{vanderweij2024}
found workable where a single summed vector fails; evaluating that composition
with both of this pipeline's instruments is a natural extension.

%=======================================================================
\section{Conclusion}

We built one pipeline that explores, for three safety behaviours in a single
model, the main ways of probing and steering: two estimators crossed with three
pooling schemes, layers selected by predictive and by causal criteria, and every
direction evaluated as both a detector and a steering vector, with the
surrounding choices fixed so the comparisons are attributable. The headline
comparison --- whether the direction that best detects a behaviour is the one
that best steers it --- resolves differently per behaviour: for deception and
high-stakes logistic regression reads best and mass-mean writes best, as the
literature implies; for sycophancy it inverts completely. Along the way, pooling
proved capable of mattering more than the estimator, the causal site sat below
the readable one for two behaviours, and the coherence axis showed that
mass-mean's causal advantage is partly a matter of steering harder rather than
better. We regard the sycophancy inversion, the scope rule --- steer the item,
never the apparatus --- and the conditions under which the causal metric is
trustworthy as the most useful contributions.

%=======================================================================
\section*{Acknowledgment}

I thank my supervisor, Professor Massimo Poesio, for his guidance
throughout. The pipeline builds on
the ARENA 3.0 interpretability curriculum \citep{arena2024} and on datasets
released by \citet{goldowsky2025}, \citet{mckenzie2025} and \citet{cincynlp2025}.
Use of generative AI tools is declared in Appendix~\ref{app:genai}.

%=======================================================================
\begin{thebibliography}{99}
\setlength{\itemsep}{2pt}

\bibitem[Alain and Bengio(2018)]{alain2018}
Alain, G. and Bengio, Y. (2018) `Understanding intermediate layers using linear
classifier probes', \textit{arXiv preprint} arXiv:1610.01644.

\bibitem[Azaria and Mitchell(2023)]{azaria2023}
Azaria, A. and Mitchell, T. (2023) `The internal state of an LLM knows when it's
lying', in \textit{Findings of EMNLP 2023}. arXiv:2304.13734.

\bibitem[Bai et al.(2022)]{bai2022}
Bai, Y., Jones, A., Ndousse, K., et al. (2022) `Training a helpful and harmless
assistant with reinforcement learning from human feedback', \textit{arXiv
preprint} arXiv:2204.05862.

\bibitem[Benton et al.(2024)]{benton2024}
Benton, J., Wagner, M., Christiansen, E., et al. (2024) `Sabotage evaluations
for frontier models', \textit{arXiv preprint} arXiv:2410.21514.

\bibitem[Bricken et al.(2023)]{bricken2023}
Bricken, T., Templeton, A., Batson, J., et al. (2023) `Towards monosemanticity:
decomposing language models with dictionary learning',
\textit{Transformer Circuits Thread}.

\bibitem[Burns et al.(2023)]{burns2023}
Burns, C., Ye, H., Klein, D. and Steinhardt, J. (2023) `Discovering latent
knowledge in language models without supervision', in \textit{International
Conference on Learning Representations (ICLR)}.

\bibitem[Chen et al.(2024)]{chen2024}
Chen, Y., Wu, A., DePodesta, T., et al. (2024) `Designing a dashboard for
transparency and control of conversational AI', \textit{arXiv preprint}
arXiv:2406.07882.

\bibitem[CincyNLP(2025)]{cincynlp2025}
CincyNLP (2025) \textit{Disentangling sycophancy: factorial evaluation
datasets}. Available at: \url{https://github.com/cincynlp/disentangle-sycophancy}
(Accessed: 30 July 2026).

\bibitem[Cunningham et al.(2023)]{cunningham2023}
Cunningham, H., Ewart, A., Riggs, L., Huben, R. and Sharkey, L. (2023) `Sparse
autoencoders find highly interpretable features in language models',
\textit{arXiv preprint} arXiv:2309.08600.

\bibitem[Elhage et al.(2022)]{elhage2022}
Elhage, N., Hume, T., Olsson, C., et al. (2022) `Toy models of superposition',
\textit{Transformer Circuits Thread}.

\bibitem[Goldowsky-Dill et al.(2025)]{goldowsky2025}
Goldowsky-Dill, N., Chughtai, B., Heimersheim, S. and Hobbhahn, M. (2025)
`Detecting strategic deception using linear probes', \textit{arXiv preprint}
arXiv:2502.03407.

\bibitem[Kantamneni et al.(2025)]{kantamneni2025}
Kantamneni, S., Engels, J., Rajamanoharan, S. and Nanda, N. (2025) `Are sparse
autoencoders useful? A case study in sparse probing', \textit{arXiv preprint}
arXiv:2502.16681.

\bibitem[Levinstein and Herrmann(2025)]{levinstein2024}
Levinstein, B.A. and Herrmann, D.A. (2025) `Still no lie detector for language
models: probing empirical and conceptual roadblocks',
\textit{Philosophical Studies}, 182(7), pp. 1539--1565.

\bibitem[Li et al.(2023)]{li2023}
Li, K., Patel, O., Vi\'egas, F., Pfister, H. and Wattenberg, M. (2023)
`Inference-time intervention: eliciting truthful answers from a language model',
in \textit{Advances in Neural Information Processing Systems (NeurIPS)}.

\bibitem[Marks and Tegmark(2024)]{marks2024}
Marks, S. and Tegmark, M. (2024) `The geometry of truth: emergent linear
structure in large language model representations of true/false datasets', in
\textit{Conference on Language Modeling (COLM)}. arXiv:2310.06824.

\bibitem[McDougall(2024)]{arena2024}
McDougall, C. (2024) \textit{ARENA 3.0: Alignment Research Engineer
Accelerator}. Available at: \url{https://github.com/callummcdougall/ARENA_3.0}
(Accessed: 2 August 2026).

\bibitem[McKenzie et al.(2025)]{mckenzie2025}
McKenzie, A., Pawar, U., Blandfort, P., Bankes, W., Krueger, D., Lubana, E.S.
and Krasheninnikov, D. (2025) `Detecting high-stakes interactions with
activation probes', \textit{arXiv preprint} arXiv:2506.10805.

\bibitem[Meng et al.(2022)]{meng2022}
Meng, K., Bau, D., Andonian, A. and Belinkov, Y. (2022) `Locating and editing
factual associations in GPT', in \textit{Advances in Neural Information
Processing Systems (NeurIPS)}.

\bibitem[Pacchiardi et al.(2024)]{pacchiardi2024}
Pacchiardi, L., Chan, A.J., Mindermann, S., et al. (2024) `How to catch an AI
liar: lie detection in black-box LLMs by asking unrelated questions', in
\textit{International Conference on Learning Representations (ICLR)}.

\bibitem[Perez et al.(2023)]{perez2023}
Perez, E., Ringer, S., Luko\v{s}i\=ut\.e, K., et al. (2023) `Discovering language
model behaviors with model-written evaluations', in \textit{Findings of ACL 2023}.

\bibitem[Rimsky et al.(2024)]{rimsky2024}
Rimsky, N., Gabrieli, N., Schulz, J., Tong, M., Hubinger, E. and Turner, A.
(2024) `Steering Llama 2 via contrastive activation addition', in
\textit{Proceedings of ACL 2024}. arXiv:2312.06681.

\bibitem[Scheurer et al.(2023)]{scheurer2023}
Scheurer, J., Balesni, M. and Hobbhahn, M. (2023) `Large language models can
strategically deceive their users when put under pressure',
\textit{arXiv preprint} arXiv:2311.07590.

\bibitem[Sharma et al.(2023)]{sharma2023}
Sharma, M., Tong, M., Korbak, T., et al. (2023) `Towards understanding sycophancy
in language models', \textit{arXiv preprint} arXiv:2310.13548.

\bibitem[Taori et al.(2023)]{taori2023}
Taori, R., Gulrajani, I., Zhang, T., et al. (2023) \textit{Stanford Alpaca: an
instruction-following LLaMA model}. Available at:
\url{https://github.com/tatsu-lab/stanford_alpaca} (Accessed: 2 August 2026).

\bibitem[Turner et al.(2023)]{turner2023}
Turner, A.M., Thiergart, L., Udell, D., et al. (2023) `Activation addition:
steering language models without optimization', \textit{arXiv preprint}
arXiv:2308.10248.

\bibitem[van der Weij et al.(2024)]{vanderweij2024}
van der Weij, T., Poesio, M. and Schoots, N. (2024) `Extending activation
steering to broad skills and multiple behaviours', \textit{arXiv preprint}
arXiv:2403.05767.

\bibitem[Vig et al.(2020)]{vig2020}
Vig, J., Gehrmann, S., Belinkov, Y., et al. (2020) `Investigating gender bias in
language models using causal mediation analysis', in \textit{Advances in Neural
Information Processing Systems (NeurIPS)}.

\bibitem[Zou et al.(2023)]{zou2023}
Zou, A., Phan, L., Chen, S., et al. (2023) `Representation engineering: a
top-down approach to AI transparency', \textit{arXiv preprint} arXiv:2310.01405.

\end{thebibliography}

%=======================================================================
\appendices

% AUTO-GENERATED by make_appendix_tables.py -- do not edit by hand.

% Regenerate after any pipeline re-run so these cannot drift from exported_probes/.



\section{Full Detection Results}

Every probe cell against every evaluation set, as exported to \texttt{exported\_probes/*\_detection.csv}. The out-of-distribution rows are the ones the read-side claim rests on; in-distribution AUROC is near-ceiling for two of the three behaviours and correspondingly uninformative.

\begin{table}[h]
\caption{Deception --- detection. AUROC with bootstrap 95\% CI; TPR at 1\% FPR (\texttt{n/a} where a set has $<$100 negatives, so 1\% FPR is unreachable).}
\label{tab:app-det-deception}
\centering
\scriptsize
\setlength{\tabcolsep}{3pt}
\begin{tabular}{llcc}
\toprule
Probe & Evaluation set & AUROC [95\% CI] & TPR@1\%FPR \\
\midrule
MM (last) & test (ID) & 0.948 [0.91,0.98] & n/a ($<$100 neg) \\
 & AI Liar (OOD) & 0.959 [0.90,0.99] & n/a ($<$100 neg) \\
 & Roleplaying (OOD) & 0.669 [0.63,0.71] & 0.011 \\
 & Deception OOD pool & 0.696 [0.66,0.73] & 0.033 \\
\addlinespace[2pt]
MM (mean) & test (ID) & 1.000 [1.00,1.00] & n/a ($<$100 neg) \\
 & AI Liar (OOD) & 0.992 [0.97,1.00] & n/a ($<$100 neg) \\
 & Roleplaying (OOD) & 0.654 [0.61,0.69] & 0.016 \\
 & Deception OOD pool & 0.679 [0.64,0.72] & 0.018 \\
\addlinespace[2pt]
LR (last) & test (ID) & 0.996 [0.99,1.00] & n/a ($<$100 neg) \\
 & AI Liar (OOD) & 0.818 [0.70,0.93] & n/a ($<$100 neg) \\
 & Roleplaying (OOD) & 0.554 [0.51,0.60] & 0.011 \\
 & Deception OOD pool & 0.571 [0.53,0.61] & 0.023 \\
\addlinespace[2pt]
LR (mean) & test (ID) & 1.000 [1.00,1.00] & n/a ($<$100 neg) \\
 & AI Liar (OOD) & 0.974 [0.93,1.00] & n/a ($<$100 neg) \\
 & Roleplaying (OOD) & 0.706 [0.67,0.74] & 0.078 \\
 & Deception OOD pool & 0.722 [0.69,0.76] & 0.075 \\
\addlinespace[2pt]
LR-attention & test (ID) & 1.000 [1.00,1.00] & n/a ($<$100 neg) \\
 & AI Liar (OOD) & 0.989 [0.97,1.00] & n/a ($<$100 neg) \\
 & Roleplaying (OOD) & 0.657 [0.62,0.69] & 0.049 \\
 & Deception OOD pool & 0.685 [0.65,0.72] & 0.055 \\
\addlinespace[2pt]
\bottomrule
\end{tabular}
\end{table}

\begin{table}[h]
\caption{High-stakes --- detection. AUROC with bootstrap 95\% CI; TPR at 1\% FPR (\texttt{n/a} where a set has $<$100 negatives, so 1\% FPR is unreachable).}
\label{tab:app-det-highstakes}
\centering
\scriptsize
\setlength{\tabcolsep}{3pt}
\begin{tabular}{llcc}
\toprule
Probe & Evaluation set & AUROC [95\% CI] & TPR@1\%FPR \\
\midrule
MM (last) & test (ID) & 0.867 [0.84,0.89] & 0.308 \\
 & anthropic\_hh\_balanced (OOD) & 0.678 [0.66,0.70] & 0.058 \\
\addlinespace[2pt]
MM (mean) & test (ID) & 0.951 [0.94,0.96] & 0.372 \\
 & anthropic\_hh\_balanced (OOD) & 0.699 [0.68,0.72] & 0.031 \\
\addlinespace[2pt]
LR (last) & test (ID) & 0.977 [0.97,0.98] & 0.688 \\
 & anthropic\_hh\_balanced (OOD) & 0.782 [0.77,0.80] & 0.128 \\
\addlinespace[2pt]
LR (mean) & test (ID) & 0.995 [0.99,1.00] & 0.828 \\
 & anthropic\_hh\_balanced (OOD) & 0.826 [0.81,0.84] & 0.080 \\
\addlinespace[2pt]
LR-attention & test (ID) & 0.997 [0.99,1.00] & 0.948 \\
 & anthropic\_hh\_balanced (OOD) & 0.808 [0.79,0.82] & 0.080 \\
\addlinespace[2pt]
\bottomrule
\end{tabular}
\end{table}

\begin{table}[h]
\caption{Sycophancy --- detection. AUROC with bootstrap 95\% CI; TPR at 1\% FPR (\texttt{n/a} where a set has $<$100 negatives, so 1\% FPR is unreachable).}
\label{tab:app-det-sycophancy}
\centering
\scriptsize
\setlength{\tabcolsep}{3pt}
\begin{tabular}{llcc}
\toprule
Probe & Evaluation set & AUROC [95\% CI] & TPR@1\%FPR \\
\midrule
MM (last) & test (ID) & 0.976 [0.97,0.98] & 0.862 \\
 & math (OOD) & 0.879 [0.86,0.90] & 0.200 \\
\addlinespace[2pt]
MM (mean) & test (ID) & 0.928 [0.91,0.94] & 0.602 \\
 & math (OOD) & 0.802 [0.77,0.83] & 0.018 \\
\addlinespace[2pt]
LR (last) & test (ID) & 0.975 [0.97,0.98] & 0.859 \\
 & math (OOD) & 0.847 [0.82,0.87] & 0.052 \\
\addlinespace[2pt]
LR (mean) & test (ID) & 0.970 [0.96,0.98] & 0.790 \\
 & math (OOD) & 0.825 [0.80,0.85] & 0.046 \\
\addlinespace[2pt]
LR-attention & test (ID) & 0.974 [0.96,0.98] & 0.865 \\
 & math (OOD) & 0.861 [0.84,0.88] & 0.050 \\
\addlinespace[2pt]
\bottomrule
\end{tabular}
\end{table}

\section{Full Steering Response Curves}

All seven coefficients for every direction and evaluation set, as exported to \texttt{exported\_probes/*\_steering\_full.json}. Three things are easiest to see here rather than in the body: whether a curve is monotone in the coefficient, whether the negative half behaves as a control, and where judge coherence falls away from its unsteered baseline.

\begin{table}[h]
\caption{Deception, ID: steering response. Readout NIE, judge behaviour and judge coherence (1--10; \texttt{---} where the judge arm is deliberately skipped).}
\label{tab:app-steer-deception-ID}
\centering
\scriptsize
\setlength{\tabcolsep}{3pt}
\begin{tabular}{llrrrrrrr}
\toprule
Direction & & -2 & -1 & -0.5 & +0 & +0.5 & +1 & +2 \\
\midrule
MM & NIE & -0.58 & -0.58 & -0.56 & +0.00 & +2.61 & +7.25 & +14.63 \\
 & judge beh. & --- & --- & --- & --- & --- & --- & --- \\
 & judge coh. & --- & --- & --- & --- & --- & --- & --- \\
\addlinespace[2pt]
LR & NIE & -0.56 & -0.55 & -0.44 & +0.00 & +0.99 & +2.29 & +3.70 \\
 & judge beh. & --- & --- & --- & --- & --- & --- & --- \\
 & judge coh. & --- & --- & --- & --- & --- & --- & --- \\
\addlinespace[2pt]
LR-attention & NIE & -0.53 & -0.49 & -0.33 & +0.00 & +0.54 & +1.15 & +2.57 \\
 & judge beh. & --- & --- & --- & --- & --- & --- & --- \\
 & judge coh. & --- & --- & --- & --- & --- & --- & --- \\
\addlinespace[2pt]
\bottomrule
\end{tabular}
\end{table}

\begin{table}[h]
\caption{Deception, AI Liar (OOD): steering response. Readout NIE, judge behaviour and judge coherence (1--10; \texttt{---} where the judge arm is deliberately skipped).}
\label{tab:app-steer-deception-AILiarOOD}
\centering
\scriptsize
\setlength{\tabcolsep}{3pt}
\begin{tabular}{llrrrrrrr}
\toprule
Direction & & -2 & -1 & -0.5 & +0 & +0.5 & +1 & +2 \\
\midrule
MM & NIE & +0.26 & -0.10 & -0.09 & +0.00 & +0.09 & +0.16 & +0.24 \\
 & judge beh. & 3.33 & 2.08 & 2.38 & 2.17 & 2.33 & 2.71 & 3.25 \\
 & judge coh. & 2.12 & 5.33 & 5.58 & 4.38 & 7.42 & 7.58 & 1.92 \\
\addlinespace[2pt]
LR & NIE & -0.37 & -0.14 & -0.07 & +0.00 & +0.05 & +0.05 & -0.15 \\
 & judge beh. & 2.38 & 2.12 & 2.21 & 2.17 & 2.42 & 2.33 & 3.12 \\
 & judge coh. & 4.96 & 5.17 & 5.58 & 4.38 & 5.58 & 5.88 & 3.83 \\
\addlinespace[2pt]
LR-attention & NIE & -0.13 & -0.04 & -0.00 & +0.00 & +0.02 & -0.02 & -0.15 \\
 & judge beh. & 2.04 & 2.67 & 2.12 & 2.17 & 2.04 & 2.33 & 2.50 \\
 & judge coh. & 4.75 & 4.79 & 5.38 & 4.38 & 6.04 & 5.54 & 6.29 \\
\addlinespace[2pt]
\bottomrule
\end{tabular}
\end{table}

\begin{table}[h]
\caption{Deception, Roleplaying (OOD): steering response. Readout NIE, judge behaviour and judge coherence (1--10; \texttt{---} where the judge arm is deliberately skipped).}
\label{tab:app-steer-deception-RoleplayingOOD}
\centering
\scriptsize
\setlength{\tabcolsep}{3pt}
\begin{tabular}{llrrrrrrr}
\toprule
Direction & & -2 & -1 & -0.5 & +0 & +0.5 & +1 & +2 \\
\midrule
MM & NIE & +0.53 & -0.02 & -0.00 & +0.00 & -0.05 & -0.07 & -0.15 \\
 & judge beh. & 6.42 & 7.12 & 6.62 & 6.96 & 6.75 & 6.62 & 6.08 \\
 & judge coh. & 1.04 & 5.00 & 5.62 & 6.08 & 6.04 & 5.71 & 1.38 \\
\addlinespace[2pt]
LR & NIE & -0.27 & -0.09 & -0.03 & +0.00 & -0.01 & -0.08 & -0.53 \\
 & judge beh. & 6.79 & 6.46 & 7.12 & 6.96 & 7.33 & 7.83 & 7.12 \\
 & judge coh. & 5.08 & 6.08 & 5.38 & 6.08 & 5.75 & 5.67 & 2.58 \\
\addlinespace[2pt]
LR-attention & NIE & -0.02 & +0.02 & +0.02 & +0.00 & -0.06 & -0.18 & -0.49 \\
 & judge beh. & 7.12 & 6.42 & 6.79 & 6.96 & 7.33 & 7.12 & 7.08 \\
 & judge coh. & 5.29 & 6.12 & 5.79 & 6.08 & 5.67 & 5.79 & 4.92 \\
\addlinespace[2pt]
\bottomrule
\end{tabular}
\end{table}

\begin{table}[h]
\caption{High-stakes, ID: steering response. Readout NIE, judge behaviour and judge coherence (1--10; \texttt{---} where the judge arm is deliberately skipped).}
\label{tab:app-steer-highstakes-ID}
\centering
\scriptsize
\setlength{\tabcolsep}{3pt}
\begin{tabular}{llrrrrrrr}
\toprule
Direction & & -2 & -1 & -0.5 & +0 & +0.5 & +1 & +2 \\
\midrule
MM & NIE & -3.06 & -3.27 & -2.01 & +0.00 & +1.25 & +1.52 & +1.24 \\
 & judge beh. & 2.46 & 3.42 & 3.71 & 3.75 & 5.33 & 5.62 & 5.62 \\
 & judge coh. & 5.04 & 4.75 & 4.46 & 4.42 & 5.00 & 4.75 & 2.83 \\
\addlinespace[2pt]
LR & NIE & -2.49 & -1.41 & -0.72 & +0.00 & +0.52 & +1.01 & +1.49 \\
 & judge beh. & 3.54 & 3.71 & 4.46 & 3.75 & 4.50 & 4.54 & 4.42 \\
 & judge coh. & 5.00 & 5.12 & 4.75 & 4.42 & 4.50 & 4.58 & 4.38 \\
\addlinespace[2pt]
LR-attention & NIE & -3.62 & -2.32 & -1.30 & +0.00 & +0.84 & +1.38 & +1.56 \\
 & judge beh. & 3.50 & 3.71 & 3.50 & 3.75 & 4.50 & 4.38 & 6.00 \\
 & judge coh. & 5.33 & 4.71 & 4.75 & 4.42 & 4.71 & 4.79 & 5.04 \\
\addlinespace[2pt]
\bottomrule
\end{tabular}
\end{table}

\begin{table}[h]
\caption{High-stakes, Anthropic HH (OOD): steering response. Readout NIE, judge behaviour and judge coherence (1--10; \texttt{---} where the judge arm is deliberately skipped).}
\label{tab:app-steer-highstakes-AnthropicHHOOD}
\centering
\scriptsize
\setlength{\tabcolsep}{3pt}
\begin{tabular}{llrrrrrrr}
\toprule
Direction & & -2 & -1 & -0.5 & +0 & +0.5 & +1 & +2 \\
\midrule
MM & NIE & -0.06 & -0.07 & -0.06 & +0.00 & +0.16 & +0.56 & +1.25 \\
 & judge beh. & 1.42 & 1.67 & 1.83 & 1.79 & 1.83 & 1.96 & 2.58 \\
 & judge coh. & 6.71 & 7.38 & 6.54 & 6.21 & 6.21 & 5.00 & 2.83 \\
\addlinespace[2pt]
LR & NIE & -0.07 & -0.05 & -0.03 & +0.00 & +0.05 & +0.12 & +0.36 \\
 & judge beh. & 1.83 & 1.79 & 1.79 & 1.79 & 1.75 & 1.71 & 1.79 \\
 & judge coh. & 5.92 & 6.29 & 6.38 & 6.21 & 6.29 & 6.21 & 5.79 \\
\addlinespace[2pt]
LR-attention & NIE & -0.07 & -0.07 & -0.04 & +0.00 & +0.09 & +0.26 & +0.78 \\
 & judge beh. & 1.79 & 1.83 & 1.83 & 1.79 & 1.75 & 1.79 & 1.96 \\
 & judge coh. & 7.00 & 6.50 & 6.54 & 6.21 & 6.00 & 5.88 & 6.00 \\
\addlinespace[2pt]
\bottomrule
\end{tabular}
\end{table}

\begin{table}[h]
\caption{Sycophancy, ID: steering response. Readout NIE, judge behaviour and judge coherence (1--10; \texttt{---} where the judge arm is deliberately skipped).}
\label{tab:app-steer-sycophancy-ID}
\centering
\scriptsize
\setlength{\tabcolsep}{3pt}
\begin{tabular}{llrrrrrrr}
\toprule
Direction & & -2 & -1 & -0.5 & +0 & +0.5 & +1 & +2 \\
\midrule
MM & NIE & +0.03 & -0.55 & +0.04 & +0.00 & +0.89 & +0.08 & +0.03 \\
 & judge beh. & 1.00 & 1.04 & 1.83 & 1.17 & 1.29 & 5.67 & 1.00 \\
 & judge coh. & 1.00 & 7.25 & 8.12 & 9.08 & 9.71 & 3.83 & 1.00 \\
\addlinespace[2pt]
LR & NIE & -0.07 & -0.04 & -0.04 & +0.00 & +0.01 & -0.00 & +0.03 \\
 & judge beh. & 1.46 & 1.50 & 1.42 & 1.17 & 1.42 & 1.42 & 1.08 \\
 & judge coh. & 9.21 & 9.29 & 9.42 & 9.08 & 9.54 & 9.00 & 9.17 \\
\addlinespace[2pt]
LR-attention & NIE & -0.32 & -0.22 & -0.14 & +0.00 & +0.11 & +0.25 & +0.51 \\
 & judge beh. & 1.38 & 3.17 & 1.46 & 1.17 & 1.04 & 1.04 & 1.42 \\
 & judge coh. & 9.21 & 9.42 & 9.29 & 9.08 & 9.33 & 8.67 & 8.33 \\
\addlinespace[2pt]
\bottomrule
\end{tabular}
\end{table}

\begin{table}[h]
\caption{Sycophancy, Math (OOD): steering response. Readout NIE, judge behaviour and judge coherence (1--10; \texttt{---} where the judge arm is deliberately skipped).}
\label{tab:app-steer-sycophancy-MathOOD}
\centering
\scriptsize
\setlength{\tabcolsep}{3pt}
\begin{tabular}{llrrrrrrr}
\toprule
Direction & & -2 & -1 & -0.5 & +0 & +0.5 & +1 & +2 \\
\midrule
MM & NIE & -0.21 & -1.05 & -0.09 & +0.00 & +0.91 & -0.16 & -0.21 \\
 & judge beh. & 1.00 & 1.00 & 1.00 & 1.00 & 3.71 & 7.33 & 1.00 \\
 & judge coh. & 1.00 & 3.46 & 6.83 & 6.58 & 9.79 & 7.33 & 1.00 \\
\addlinespace[2pt]
LR & NIE & -0.05 & -0.06 & -0.03 & +0.00 & -0.03 & -0.01 & -0.02 \\
 & judge beh. & 1.75 & 1.00 & 1.00 & 1.00 & 1.00 & 1.04 & 1.00 \\
 & judge coh. & 7.25 & 6.79 & 6.29 & 6.58 & 6.67 & 6.21 & 7.17 \\
\addlinespace[2pt]
LR-attention & NIE & -1.03 & -0.61 & -0.19 & +0.00 & +0.13 & +0.34 & +0.48 \\
 & judge beh. & 1.38 & 2.50 & 1.83 & 1.00 & 1.00 & 1.04 & 1.04 \\
 & judge coh. & 7.88 & 8.25 & 8.12 & 6.58 & 6.92 & 7.00 & 7.25 \\
\addlinespace[2pt]
\bottomrule
\end{tabular}
\end{table}


\section{Reproducibility}
\label{app:repro}
Model \texttt{meta-llama/Llama-3.1-8B-Instruct} in 8-bit, seed 42, on one NVIDIA RTX 5090. Selected
layers and blocks are in Table~\ref{tab:config}. Activations are extracted once per example at all 32
layers and cached as fp16; the cache key is content-blind, so a loader change requires deleting the
matching cache file.

\emph{Provenance and exact inputs.} Every prompt, token mask and per-example token count is documented
in the repository's data reference, where each block shown is a real pull from the seed-42 splits
rather than a reconstruction. Package docstrings cite the source-paper section for each estimator, so
a claim in §III can be traced to the line that implements it.

\emph{Splits.} Pair-aware throughout: by fact for deception, by pair for high-stakes, by prompt group
for sycophancy. Layer and block selection use training data only; the block-selection sweep uses
train-split transcripts while every reported curve uses held-out ones.

\emph{Regenerating the appendix.} The detection tables are generated from
\texttt{exported\_probes/} by \texttt{make\_appendix\_tables.py}, so re-running the pipeline and
re-running that script is sufficient to bring them back into agreement with the artefacts.

% The remaining appendices are figure-dominated; single-column layout lets each
% figure take its natural width and keeps headings, figures and captions together.
% \raggedbottom stops \flushbottom stretching short pages into sparse ones.
\onecolumn
\raggedbottom

\section{Detection-Layer Selection Curves}\label{app:layer-curves}

The body shows one behaviour per diagnostic to stay within the page limit. This appendix and the
following ones supply the other behaviours, read side then write side, so a reader checking whether a
claim generalises does not have to take it on trust. Every caption states what the figure shows, why
it is included, and what to look for.

\begin{figure}[H]
\centering
\figslot{deception/layers.png}{0.55\textwidth}{1.8cm}\\[4pt]
\figslot{high_stakes/layers.png}{0.55\textwidth}{1.8cm}
\caption{\textbf{Deception (top) and high-stakes (bottom).} Pair-aware 5-fold CV
AUROC against layer, mass-mean as selector with a logistic-regression overlay as a robustness check,
chosen layer marked; counterpart to Fig.~\ref{fig:layers}, which shows sycophancy.
\emph{Why included.} The three curves show the range of shapes the same procedure
produces, and that the LR overlay tracks the MM selector: if it did not, the layer would be an
artefact of the cheap estimator used to choose it. \emph{What to look for.} High-stakes has a genuine
if broad peak (0.949 at layer 20), so its argmax is meaningful. Deception saturates near 1.000 from
roughly layer 11 onward (LR from layer 1), so its argmax at 12 is weakly determined: the reported
detection numbers would barely move anywhere in that band. Sycophancy (Fig.~\ref{fig:layers}) rises
late and peaks at 19, unusually deep relative to both other behaviours and to the layers reported by
Apollo (22 of 80) and McKenzie et al.\ (31 of 80); its logistic-regression overlay is already at 0.97
by layer 8, so the \emph{late} layer is a property of the mass-mean representation, not of the
behaviour.}
\label{fig:app-layers}
\end{figure}

\section{Linear Structure at the Detection Layer}\label{app:linear}

\begin{figure}[H]
\centering
\figslot{deception/pca.png}{0.72\textwidth}{2.0cm}
\caption{\textbf{Linear structure, deception, at the detection layer.} \emph{Left:} PCA of centred
activations, showing no class separation (PC1 carries 8.3\% of variance). \emph{Right:} projections
onto the trained mass-mean direction, showing clean bimodal separation. \emph{Why both panels.}
Separation under unsupervised PCA would suggest the probe rides a high-variance topic confound of the
kind Apollo report; separation only along the supervised direction argues the confound axes are
roughly orthogonal to it. \emph{What to look for.} The co-located point pairs on the left are the two
members of each instructed pair, byte-identical apart from the system prompt: direct visual
confirmation that the behaviour is a low-variance direction that content variance swamps.}
\label{fig:pca}
\end{figure}

\begin{figure}[H]
\centering
\figslot{high_stakes/pca.png}{0.72\textwidth}{1.6cm}\\[4pt]
\figslot{sycophancy/pca.png}{0.72\textwidth}{1.6cm}
\caption{\textbf{High-stakes (top) and sycophancy (bottom).} PCA of centred activations (left of each
pair) beside projections onto the trained direction (right). \emph{Why included.} To show the
PCA-flat/projection-separated pattern is not specific to deception. \emph{What to look for.} Both
scatters show pronounced clusters, but the clusters are \emph{mixed-label}: they track dataset source
for high-stakes and claim domain for sycophancy, not the class. High-stakes is the case where label
separation in PCA \emph{would} have been unsurprising, since its classes differ in content rather
than only in framing; that the label stays unseparated is the stronger version of the argument, and
is why the projection panel rather than the scatter is treated as load-bearing throughout. Sycophancy's
projection histogram is additionally bimodal \emph{within} the positive class --- one mode overlapping
the negatives and one far mode --- reflecting the domain structure of the factorial data.}
\label{fig:app-pca}
\end{figure}

\section{Probe-Direction Cosines}\label{app:cosines}

\begin{figure}[H]
\centering
\figslot{deception/cosines.png}{0.32\textwidth}{1.8cm}\hfill
\figslot{high_stakes/cosines.png}{0.32\textwidth}{1.8cm}\hfill
\figslot{sycophancy/cosines.png}{0.32\textwidth}{1.8cm}
\caption{\textbf{Deception (left), high-stakes (centre), sycophancy (right):} pairwise cosine
similarity between the five probe directions at each behaviour's detection layer. \emph{Why
included.} This addresses the obvious objection to the whole comparison, that the estimators might be
finding the \emph{same} direction and differing only in scale, in which case ``which steers better''
would be a units artefact; cosines well below 1 mean they are genuinely different directions.
\emph{What to look for.} Alignment falls with each behaviour. Deception's directions are the most
aligned, and its attention probe's direction is closer to the mass-mean family (0.74--0.81) than to
the LR family (0.54--0.55) despite sharing LR's training objective. High-stakes' cross-family cosines
sit at 0.26--0.53. Sycophancy's directions are the least mutually aligned of the three --- its
mass-mean and logistic-regression directions are near-orthogonal (0.07--0.13), and even its two
mass-mean poolings agree only at 0.46 --- and it is the informative case: its inversion in §IV-C is
\emph{not} two copies of one direction behaving differently, but genuinely different directions, of
which all five read the behaviour and only some write it.}
\label{fig:app-cos}
\end{figure}

\section{Attention-Probe Weights}\label{app:attnw}

\begin{figure}[H]
\centering
\figslot{sycophancy/attn_sycophantic.png}{0.95\textwidth}{2.0cm}
\caption{\textbf{Attention-probe weights, sycophancy.} A softmax over the masked positions only, so
weights sum to one within the mask and are comparable across examples. \emph{Why included.} The
attention direction is sycophancy's \emph{best} steerer and the only one with a monotone steering
response, so where it places mass is evidence about what the writable feature actually is. It is also
the closest thing this pipeline has to a localisation result, activation patching being unavailable on
dialogues (§III-D). \emph{What to look for.} The mass sits almost entirely on the answer: on this
example the echoed country token carries the maximum weight (0.38), with the trailing period next,
while the templated ``I believe the city of \ldots{} is in'' prefix receives near-uniform low weight
--- the probe has learned to ignore precisely the span that dilutes mean pooling (§III-B). The
contrast with high-stakes is sharp: there the learned weights are diffuse (maximum 0.03 on any
token), matching a label that is a property of the whole situation rather than of one word.}
\label{fig:app-attn}
\end{figure}

\section{Per-Token Probe Scores}\label{app:tokens}

\begin{figure}[H]
\centering
\figslot{deception/heatmap_deceptive.png}{0.9\textwidth}{1.6cm}\\[3pt]
\figslot{deception/heatmap_honest.png}{0.9\textwidth}{1.6cm}
\caption{\textbf{Per-token probe scores, one deception pair:} deceptive persona (top), honest twin
(bottom) --- the mask of §III-B made visible, since only assistant-content tokens are scored.
\emph{Why a pair rather than one example.} The two members share an identical user turn and the same
truncated fact, differing only in the system persona, so any difference in colour is the probe reading
\emph{context} rather than lexical content. \emph{What to look for.} Where the score rises: early,
before any giveaway, is the intent reading Apollo report; only on a tell-tale phrase would suggest a
surface cue.}
\label{fig:heat}
\end{figure}

\section{Per-Layer NIE Sweeps}\label{app:niesweeps}

\begin{figure}[H]
\centering
\figslot{deception/niesweep.png}{0.32\textwidth}{1.8cm}\hfill
\figslot{high_stakes/niesweep.png}{0.32\textwidth}{1.8cm}\hfill
\figslot{sycophancy/niesweep.png}{0.32\textwidth}{1.8cm}
\caption{\textbf{Deception (left), high-stakes (centre), sycophancy (right).} Per-layer NIE with one
layer steered at a time; the selected block is the contiguous five-layer window of maximum summed
effect. \emph{Why included.} The three together show whether the read/write layer offset is a property
of the pipeline or of the behaviour. \emph{What to look for.} Deception's causal band (11--15)
encloses its detection layer (12) --- the one behaviour where the two criteria agree on a region ---
though its effect is concentrated in two sharp spikes at layers 11 and 13, with layer 12 itself
slightly negative. High-stakes' band (6--10) and sycophancy's (6--10) both sit far below their own
detection layers (20, 19). High-stakes's profile is a clean unimodal bump an order of magnitude larger
than the other two (peak 1.35 at layer 9, zero above 14). Note what this rules out: the offset does
\emph{not} track the dissociation inversion --- high-stakes has the largest offset yet its
dissociation holds, while sycophancy, with the same offset, inverts. Layer geometry alone does not
explain which behaviours dissociate.}
\label{fig:app-nie}
\end{figure}

\section{In-Distribution Steering}\label{app:steer-id}

\begin{figure}[H]
\centering
\figslot{deception/steering_id.png}{0.56\textwidth}{1.6cm}\\[1pt]
{\footnotesize (a) \textbf{deception}, ID: judge columns empty by design}\\[5pt]
\figslot{high_stakes/steering_id.png}{0.56\textwidth}{1.6cm}\\[1pt]
{\footnotesize (b) \textbf{high-stakes}, ID}\\[5pt]
\figslot{sycophancy/steering_id.png}{0.56\textwidth}{1.6cm}\\[1pt]
{\footnotesize (c) \textbf{sycophancy}, ID}
\caption{\textbf{Steering, in-distribution.} Panels as Fig.~\ref{fig:steer}: readout NIE, judge
behaviour, judge coherence. \emph{Why in the appendix rather than the body.} Two independent reasons.
Steering blocks were selected to maximise NIE on training transcripts, so these values are closer to
fitted than held-out. And deception's in-distribution readout is degenerate: the model answers ``NO''
to the deception question for \emph{both} classes ($PD^{+}=-0.906$, $PD^{-}=-0.964$), leaving a class
gap of $0.058$, so its NIE should not be quoted. Its judge columns are empty for a separate reason,
the in-distribution generation prompts are byte-identical under greedy decoding, so the judge would
score one generation repeatedly. \emph{What to look for.} Whether the direction \emph{ordering}
matches the out-of-distribution panels; the magnitudes here should not be compared with them.}
\label{fig:app-steer-id}
\end{figure}

\section{Deception Steering on Roleplaying}\label{app:steer-rp}

\begin{figure}[H]
\centering
\figslot{deception/steering_roleplaying.png}{0.78\textwidth}{2.0cm}
\caption{\textbf{Steering deception on roleplaying, the second out-of-distribution set.} \emph{Why
included.} Deception is the only behaviour with two out-of-distribution sets, deliberately of different
kinds: AI Liar is context-defined, with byte-identical answers across labels, while roleplaying is
content-defined, with genuinely different completions. Reporting only the former would overstate
transfer. \emph{What to look for.} Transfer is weak but not absent: mass-mean reaches NIE $+0.13$ at
coefficient $1$ and $+0.35$ at $2$, the attention direction $+0.21$ at $2$, with plain logistic
regression flat --- a faint echo of the detection collapse on the same set
(Table~\ref{tab:app-det-deception}) rather than a full mirror of it. Mass-mean's curve is also U-shaped, rising at \emph{both} sweep ends;
its $-2$ point coincides with a coherence drop to 2.96, so the negative-end rise is degradation, not
signal. The judge panel is separately uninformative: roleplaying prompts have no behaviour-off
variant, since \texttt{roleplaying\_messages} varies only the assistant completion and the generation
prompt discards it, and the unsteered baseline already scores 6.71/10 deceptive, so there is little
headroom to induce into (the judged scores span 5.58--7.96 across all directions and coefficients).}
\label{fig:app-steer-rp}
\end{figure}

\section{The Interface in Use}
\label{app:ui-shots}

Figs.~\ref{fig:ui} and \ref{fig:ui-steer} show the detection and steering tabs; these examples show
the multi-behaviour mode. Both are live generations on free input, far from any evaluation set, with the same calibrated
directions the results sections report, each behaviour's directions applied at its own steering
block.

\begin{figure}[H]
\centering
\begin{minipage}[t]{0.485\textwidth}
\centering
\figslot{ui/multi-steering/deception-sycophancy.png}{\linewidth}{2.2cm}
\caption{\textbf{Two behaviours composed (deception $+1.0$, sycophancy $+0.75$,
both mass-mean),} each direction applied at its own behaviour's steering block,
following \citet{vanderweij2024}. Unsteered, the model declines to endorse
swimming the Channel without a life vest; steered, it caves into stumbling
enthusiasm for the dangerous plan --- the sycophantic caving of §IV-C composed
with deception, on out-of-distribution input. The in-app banner marks the mode
exploratory: how directions compose is an open question the pipeline does not
yet evaluate.}
\label{fig:ui-multi}
\end{minipage}\hfill
\begin{minipage}[t]{0.485\textwidth}
\centering
\figslot{ui/multi-steering/deception-high-stakes.png}{\linewidth}{2.2cm}
\caption{\textbf{Composition can also fail, informatively (deception $+1.75$,
high-stakes $+1.5$).} At higher combined coefficients the steered generation
collapses into alarmist non-sequitur (``I can't help but say that the sky is
falling.'') where the baseline suggests cuisines --- the degradation the
coherence axis quantifies in §IV-D, reproduced interactively. The two
multiplier sliders make the coherence cliff something a user can find by
feel.}
\label{fig:ui-multi2}
\end{minipage}
\end{figure}

\section{The Steering-Scope Ablation}
\label{app:scope}

Three scopes for the generation-side intervention were compared on otherwise
identical runs, after judged dose-response curves flattened when the scope was
first narrowed to match published practice:

\begin{enumerate}
  \item \emph{Completion-only}: steer generated tokens only. The first generated
  token's logits are computed from the hidden state at the final \emph{prompt}
  position, which this scope never touches, so the reply's opening is chosen
  unsteered.
  \item \emph{CAA-exact}: steer inclusively from the last instruction-end token,
  reproducing the published CAA code path \citep{rimsky2024}. This repairs the
  first-token gap but still leaves the item's representation unmoved.
  \item \emph{Whole-item}: steer every prompt position and every generated token
  (the convention adopted, Fig.~\ref{fig:scope}).
\end{enumerate}

Under both narrower scopes, judged behaviour stayed within noise of its unsteered
baseline at every coefficient for all three behaviours, while whole-item steering
produces the graded dose-responses of §\ref{sec:results}. The readout NIE is
byte-identical across the three regimes, as it must be: its scope (the item
transcript within the few-shot prompt) is fixed separately and no generation is
involved. The archived per-regime curves are in the repository under
\texttt{research-report/figures/scope\_ablation/}.

The reading is not that CAA's published results are wrong, but that its scope is
not sufficient \emph{here}: CAA steers a 7-token answer-selection context where
the behaviour is carried by the immediate decision, whereas these generation
prompts embed the behaviour in the item's representation (a scenario, a
situation, a user's claim), which must itself be moved for the completion to
change. Hence the rule stated in §\ref{sec:steering}: steer the item, never the
apparatus.

\section{Statement on the Use of Generative AI Tools}
\label{app:genai}
The following is a factual record of generative-AI use, for transfer into the
departmental declaration template.
\emph{(i) Code.} An AI coding assistant was used to help implement and refactor the package, write
driver-notebook scaffolding, and diagnose defects (for example the double-BOS tokenisation defect in
§III-G). All code was reviewed and executed by the author, and all reported numbers were produced by
running it.
\emph{(ii) Literature comparison.} An AI assistant was used to cross-check the implementation against
the source papers section by section; the author verified each discrepancy against the papers
directly.
\emph{(iii) Writing.} An AI assistant was used to draft and edit prose. The research questions,
experimental design, results and conclusions are the author's own.
\emph{(iv) As an experimental component.} Claude Haiku 4.5 is used \emph{within} the methodology as
the LLM judge scoring steered generations (§III-F); this is a measured component of the system, not a
writing aid.

\end{document}
